\documentclass[11pt]{article}
\usepackage[margin=1in]{geometry}
\usepackage{amsmath,amssymb,amsthm,mathtools,microtype}
\usepackage[hidelinks]{hyperref}
\hypersetup{
 pdftitle={The planar matching endpoints and their leading constant},
 pdfauthor={Michael Goldman, Francesco Grotto, Dario Trevisan}}
\newenvironment{aiuse}
  {\renewcommand{\abstractname}{AI use}\begin{abstract}}
  {\end{abstract}}
\newtheorem{theorem}{Theorem}
\newtheorem{lemma}{Lemma}
\newtheorem{proposition}{Proposition}
\newtheorem{corollary}{Corollary}
\newtheorem{remark}{Remark}
\newcommand{\abs}[1]{\lvert#1\rvert}
\newcommand{\norm}[1]{\lVert#1\rVert}
\newcommand{\one}{\mathbf{1}}
\newcommand{\1}{\mathbf{1}}
\newcommand{\supp}{\operatorname{supp}}
\newcommand{\op}{\mathrm{op}}
\newcommand{\diam}{\operatorname{diam}}
\newcommand{\Lip}{\operatorname{Lip}}
\newcommand{\Cov}{\operatorname{Cov}}
\newcommand{\Id}{\operatorname{Id}}
\newcommand{\cF}{\mathcal F}
\newcommand{\cN}{\mathcal N}
\newcommand{\less}{\lesssim}
\DeclareMathOperator{\arsinh}{arsinh}
\DeclareMathOperator{\diagop}{diag}
\DeclareMathOperator{\per}{Per}
\allowdisplaybreaks[1]
\numberwithin{equation}{section}
\newcommand{\E}{\mathbb E}
\newcommand{\PP}{\mathbb P}
\newcommand{\R}{\mathbb R}
\newcommand{\T}{\mathsf T}
\newcommand{\dist}{\operatorname{dist}}
\newcommand{\Div}{\operatorname{div}}
\newcommand{\Var}{\operatorname{Var}}
\newcommand{\sgn}{\operatorname{sgn}}
\newcommand{\av}{\operatorname{av}}
\newcommand{\tr}{\operatorname{tr}}
\newcommand{\diag}{\operatorname{diag}}
\title{The planar matching endpoints and their leading constant}
\author{Michael Goldman \and Francesco Grotto \and Dario Trevisan}
\date{}
\begin{document}
\maketitle

\begin{aiuse}
This document was generated by an LLM through interaction with the authors,
who do not guarantee its correctness. It is made publicly available through
the GitHub repository~\cite{umc:repository}, together with a partially
revised version of the arguments, in the interest of transparency regarding
its development process. The authors plan to provide a fully digested
version of the arguments in the near future.
\end{aiuse}

\begin{abstract}
For two independent uniform samples of size \(N\) in the unit
square, the minimum over permutations of the maximum of the
horizontal quadratic Luxemburg norm and the vertical supremum
has leading constant \(1/\sqrt{2\pi}\) at scale
\(\sqrt{\log N/N}(\log\log\log N)^{1/4}\), in probability.
The same holds with the axes exchanged. We prove matching
Gaussian divergence bounds, using a state-dependent local-time clock
for the lower bound and a moving phase with a sharp covariance estimate
for the upper bound. We prove the local estimates for both constructions
and transfer the Gaussian constant to the point matching problem.
\end{abstract}

\section{Statement}

Let \(X_1,\ldots,X_N\) and \(Y_1,\ldots,Y_N\) be independent
uniform samples in \(\Omega=[0,1]^2\).
For a permutation \(\pi\), write
\(\Delta_i(j)=X_j^{(i)}-Y_{\pi(j)}^{(i)}\).
Use the Young function \(\Theta(t)=e^{t^2}-1\) and the
Luxemburg convention that the modular at the chosen scale
is at most one, with mass \(1/N\) per pair. Put
\[
 M_N=\min_\pi\max\{\|\Delta_1\|_{\psi_2},
                         \|\Delta_2\|_\infty\}.
\]

\begin{theorem}\label{thm:main}
With natural logarithms,
\[
 \frac{M_N}{\sqrt{\log N/N}\,(\log\log\log N)^{1/4}}
 \xrightarrow{\mathbb P}\frac{1}{\sqrt{2\pi}} .
\]
The same conclusion holds for
\(\min_\pi\max\{\|\Delta_1\|_\infty,\|\Delta_2\|_{\psi_2}\}\).
\end{theorem}

The Gaussian result uses continuum currents with an exact
zero-normal divergence condition. For centered white noise
\(W^0\) on \(\Omega\), write \(P_hW^0\) for its dyadic histogram.
We use the flux gauge
\[
 \mathcal F(f)=\inf_{\operatorname{div}J=f,\ J\cdot n=0}
      \max\{\|J_1\|_{\psi_2(dx)},\|J_2\|_{L^\infty(dx)}\}.
\]
The infimum is over integrable fields on the square with weak
divergence \(f\) and zero outer normal flux. Thus
\(\int f\varphi=-\int J\cdot\nabla\varphi\) for every smooth
test on the closed square. Both norms use Lebesgue measure on
the unit square. By white noise we mean the centered Gaussian
linear process \(W(g)\), \(g\in L^2(\Omega)\), with covariance
\(\mathbb E[W(g)W(k)]=\int gk\). Its spatial centering is
\(W^0(g)=W(g-\int g)\). A histogram \(P_hW^0\) is the density
constant on the dyadic squares of side \(h\) with those white
noise masses.

\begin{theorem}\label{thm:gaussian}
As \(h\downarrow0\) through dyadic widths,
\[
 \frac{\mathcal F(P_h W^0)}
 {\sqrt{\log(1/h)}\,(\log\log\log(1/h))^{1/4}}
 \xrightarrow{\mathbb P}\frac{1}{\sqrt{2\pi}} .
\]
The expected gauge has the same asymptotic.
\end{theorem}

The normalization has two independent factors. A planar
Poisson block has coordinate covariance \(1/(4\pi)\) per
unit of natural log-scale. The determinant-one control
has optimal balanced coefficient \(\sqrt{2}\) in the
Brownian variance clock. The two samples supply a
\(\sqrt{2}\) noise factor and their spatial cutoff has
\(\log(1/h)\sim\tfrac12\log N\), which cancel in the final
normalization.


Here and below, all asymptotics are as the indicated mesh tends to zero
or the sample size tends to infinity. Constants in \(O(\cdot)\) and
\(\lesssim\) may depend on parameters explicitly fixed before this limit.

The lower construction tests the divergence against a potential with
positive vertical derivative. Size bias by this derivative makes the
ratio of the two derivatives a martingale. The clock for that martingale
determines the lower coefficient. For the upper construction, a
determinant-one chart changes the covariance of each Gaussian current.
A convex phase determines the chart and controls both coordinate costs.
We develop the local constructions and repair estimates before using
them. The last section derives the Gaussian comparison at the required
precision and proves Theorem~\ref{thm:main}.

\tableofcontents
\medskip

\section{Predictable charts and local divergence inverses}
\label{sec:local-geometry}
\label{sec:sharp-local-inverse}

A chart on a physical cell has the form \(z\mapsto Az+b\), where
\(\det A=1\). Write \(G=AA^{\mathsf T}\) for its metric. A source
\(f\) and current \(J\) in chart coordinates push forward to
\(f(A^{-1}(x-b))\) and \(AJ(A^{-1}(x-b))\), respectively. The
change of variables in the weak divergence identity proves that
this push-forward preserves divergence and zero normal flux.
If \(K_0^{-1}G\preceq G_+\preceq K_0G\), the two chart
changes have operator norm at most \(\sqrt{K_0}\).
Indeed the squared singular values of \(A_+^{-1}A\) are
the eigenvalues of \(A_+^{-1}G A_+^{-\mathsf T}\), which lie
in \([K_0^{-1},K_0]\).

\begin{lemma}[Predictable Gaussian bridges]\label{lem:predictable-bridge}
In a finite nested sequence of partitions, choose each new
partition measurably from the earlier Gaussian cell masses.
Conditional on the current history, the unrevealed noise in
each current cell is an independent white noise bridge on
that cell. In particular, once a refinement is chosen, its
new contrast has covariance equal to the orthogonal projection
from child constants onto the subspace of zero parent means.
Two successive refinements chosen together can reveal their
contrasts independently, conditional on the old history.
\end{lemma}

\begin{proof}
A bridge on \(Q\) is isonormal on
\(\{g\in L^2(Q):\int_Qg=0\}\), extended by zero on constants.
The statement starts with this decomposition on the initial
partition, or on the whole square if only the global mean is
known. Conditional on a fixed old history, the child indicators
are fixed functions in this Hilbert space after subtracting their
parent averages. Its orthogonal decomposition into the span
of these centered indicators and its orthogonal complement
gives two independent Gaussian families. The complement is
the orthogonal direct sum of the mean-zero \(L^2\) spaces of
the children. This supplies the next conditional law and proves
the claim by induction. This argument also constructs the
conditional laws for random partitions. For each realized old
history use the indicated Gaussian projection kernels on a
countable determining set of \(L^2\) tests, and then use
mean-square completion for additional tests. The geometric
partitions below have measurable indicator maps into \(L^2\)
and measurable projections, so these kernels are measurable.
A second predetermined refinement adds an orthogonal subspace
to the first and hence gives independent contrasts. An
area-preserving chart is an isometry on the test spaces.
\end{proof}

\begin{lemma}[Whitney subdivision with a minimum side]
\label{lem:local-geometry}
Let \(Q\) be a parallelogram of diameter \(D\) containing a ball
of radius \(c_0D\). Fix a minimum side \(s>0\) and a cap
\(\theta D\) with \(0<\theta<1\). For a sufficiently large
constant \(C_0\) depending only on \(c_0\), if
\(\theta D<C_0s\) we may stop without splitting. Otherwise
there is a deterministic measurable subdivision into square
children and terminal children. Assign to a square its side
\(\ell_S\), and to a terminal child the value \(\ell_S=s\).
Then
\[
 |S|\asymp\ell_S^2,\qquad
 \operatorname{diam}(S)\lesssim\ell_S,\qquad
 \ell_S\asymp\max\{s,\min(\theta D,d_Q(x))\},\quad x\in S,
 \qquad d_Q(x)=\operatorname{dist}(x,\partial Q).
\]
The constants are uniform in \(\theta\). For a sufficiently
small fixed \(c>0\), each neighborhood
\(S^+=\{x\in Q:\operatorname{dist}(x,S)<c\ell_S\}\)
meets only a bounded number of children, of comparable sides.
The neighborhoods have bounded overlap. If \(l<l'\) are two
sides occurring in \(Q\), the children of side \(l\) lie within
distance \(Cl\) of \(\partial Q\). Thus a ball of radius \(Cl'\)
can meet at most \(C' l'/l\) such children, also after enlarging
their diameters by a fixed factor.
\end{lemma}

\begin{proof}
Use the fixed dyadic lattice with smallest side \(s\), and choose
maximal lattice squares contained in \(Q\), of side
\(2s\le\ell\le\theta D\), and of boundary clearance at least
\(\gamma\ell\). Here \(\gamma\) is a large fixed constant. They
are disjoint up to null sets. The remaining intersections
with lattice \(s\)-squares cover the remaining area. To handle
a clipped \(s\)-square, take a point \(x\) in its intersection
with \(Q\) and move a fraction \(A s/D\) of the way to the
inball center. Convexity guarantees an inball of radius
\(A c_0s\) at the new point. Choose \(A\) fixed so this radius
exceeds \(3s\), and \(C_0\) large enough so the fraction is
less than one. The \(s\)-square containing this new point is
full and within \(O(s)\) of the clipped one. It is also within
\(O(s)\) of \(\partial Q\). Choose \(\gamma\) larger than this
last constant so the full square cannot belong to a selected
square. Assign each clipped piece to such an uncovered full
square, with ties broken in lattice order. Lattice packing
bounds the number of assignments per receiver. Every uncovered
full square has boundary distance \(O(s)\), because its
\(2s\)-ancestor is under the cap and fails containment or
clearance. This proves the terminal estimates, including the
area lower bound. The construction is measurable since square
containment, clearance and receiver comparisons are measurable,
and the bounded lattice search and tie rules can be fixed
in advance.

For a selected square, its dyadic parent either fails
containment or clearance, giving \(d_Q\le C\ell_S\) throughout
the child, or exceeds the cap, giving \(\ell_S>\theta D/2\).
Clearance and the cap give the converse inequality. This proves
the side comparison. For \(x\) within \(c\ell_S\) of \(S\),
the same comparison holds with \(\ell_S\) on the left. The
upper comparison follows by the Lipschitz property of \(d_Q\)
and the size floor. For the lower comparison use the clearance
when \(S\) is a square and the size floor otherwise. Intersecting
neighborhoods are comparable as well. For example, if
\(\ell_T\le\ell_S\), any point of the intersection lies within
\(c\ell_S\) of \(S\) and within \(c\ell_T\) of \(T\), so both
sides are comparable to the same clipped boundary distance.
The area bound and disjointness now give the overlap estimates.
If \(l<l'\) occurs, a selected square of side \(l\) cannot
have a parent failing the cap, since \(2l\le l'\le\theta D\).
Thus it lies in a boundary strip of width \(Cl\), as do the
terminal children. The area of this strip in a ball of radius
\(C'l'\) is \(O(l'l)\), by covering it with the four strips
along the sides of \(Q\). Dividing by the child area proves
the count. A fixed enlargement of a child only enlarges
the radius of the ball used in the count.
\end{proof}

With a fixed small cap, the nonterminal child side contracts
by an arbitrarily small fixed factor at each step. The old
square is a shape-regular parallelogram in the new chart by
the relative metric bound. Thus the comparison of its diameter
with its old side is uniform. If the subdivision uses minimum
side \(Rs\) and then a uniform \(R\times R\) split of each
nonterminal base square, the nonterminal contraction is at
most \(C/R\). In the fine split all children have their common
side and no boundary clearance is required. Base terminal cells
and unsplit parents have diameter \(O(Rs)\) in the new chart.
After a fixed-cap split, unsplit and merged terminal cells
instead have diameter \(O(s)\). All children that are retained
in the partition have area \(\gtrsim s^2\). If sides contract
by \(\sigma\), a descendant \(j\) stages later has diameter at
most \(C(\sqrt{K_0}\sigma)^j\ell_S\) in the chart of its
ancestor \(S\). This follows by composing the relative chart
bounds. It permits the ancestor oscillations controlled by
scaled H\"older seminorms to be summed geometrically by choosing
\(\sqrt{K_0}\sigma<1\).

Let \(\eta_Q\) be a smooth probability density supported in a
concentric subball of the inball, with
\(\|\eta_Q\|_\infty\lesssim D^{-2}\) and
\(\|\nabla\eta_Q\|_\infty\lesssim D^{-3}\).
The ray kernel from source \(y\) is
\[
 K_Q(x,y)=(x-y)\int_1^\infty t\,\eta_Q(y+t(x-y))\,dt .
\]
Indeed, send a unit source at \(y\) along the straight segment
to \(z\), averaged with \(\eta_Q(z)\,dz\). The weak gradient
pairing is \(\int\eta_Q(z)\phi(z)\,dz-\phi(y)\), so the
divergence is \(\delta_y-\eta_Q\) on the whole plane, with the
current extended by zero off \(Q\). The kernel formula follows
from \(x=y+(z-y)/t\). Thus the ray inverse is exact for
mean-zero sources, including the boundary condition. Direct
differentiation gives, away from \(x=y\),
\[
\begin{aligned}
 |K_Q(x,y)|&\lesssim |x-y|^{-1},\\
 |\nabla_x K_Q(x,y)|+|\nabla_y K_Q(x,y)|&\lesssim |x-y|^{-2},\\
 |(\nabla_x+\nabla_y)K_Q(x,y)|&\lesssim (D|x-y|)^{-1}.
\end{aligned}
\]
For these estimates use \(t|x-y|\le D\) in the integral.
In the last derivative the argument of \(\eta_Q\) has
simultaneous translation derivative one, and the prefactor has
derivative zero. The last estimate is valid distributionally
as well, by differentiating under simultaneous translation
with \(x-y\) fixed. There is also the support bound
\begin{equation}\label{eq:ray-support-self-contained}
 K_Q(x,y)=0\quad\text{when }|x-y|> C d_Q(x).
\end{equation}
When the integrand is nonzero, \(x=(1-t^{-1})y+t^{-1}z\)
with \(z\) in the central ball. A nearest supporting line
has distance at least \(cD\) from \(z\) and nonnegative
distance from \(y\), so \(d_Q(x)\ge cD/t\).
In particular \(\int |K_Q(x,y)|\,dy\lesssim d_Q(x)\).
The two-pole difference bound
\(\int |K_Q(x,y)-K_Q(x',y)|\,dy
 \lesssim d(1+\log(D/d))\), \(d=|x-x'|\le D/2\),
follows by integrating the singular bound within distance
\(2d\) of the points and the derivative bound farther away
along their connecting segment. This gives a continuous
representative of the current of a bounded source, vanishing
on \(\partial Q\).

\begin{lemma}[Coarse Gaussian current and envelopes]
\label{lem:coarse-envelope}
For the fixed-cap Whitney split, let \(P\) be child averaging
and \(P_Q\) parent averaging. For \(x,x'\in S^+\), the
kernels \(PK_Q(x,\cdot)\) and
\[
 \left(\frac{\ell_S}{|x-x'|}\right)^\beta
           P(K_Q(x,\cdot)-K_Q(x',\cdot)),\qquad 0<\beta<1 ,
\]
are supported in a ball of radius \(C\ell_S\) centered at a
representative of \(S\), and have supremum at most \(C_\beta/\ell_S\).
The same holds for \(\ell_S P(x,\cdot)\), the normalized density
test, where \(P(x,\cdot)=\mathbf1_T/|T|\) when \(x\in T\).
Consequently the contrast current
\(V(x)=W((P-P_Q)K_Q(x,\cdot))\) has pointwise covariance
bounded by \(CI\). If \(f=(P-P_Q)W\), then for each fixed
\(0<\alpha<1\) the even envelope
\[
 E_S=\sup_{S^+}|V|+\ell_S^\alpha[V]_{C^\alpha(S^+)}
                         +\ell_S\|f\|_{L^\infty(S^+)}
\]
has conditional tail \(\mathbb P(E_S>t\mid\mathrm{past})\le C e^{-ct^2}\).
The same conclusions hold for any fixed upper cap, with
constants depending on that cap.
\end{lemma}

\begin{proof}
For fixed cap, Lemma~\ref{lem:local-geometry} gives
\(d_Q(x)\lesssim\ell_S\) on \(S^+\). If \(S\) is a square,
\(d_Q(x)\gtrsim\ell_S\) there, by clearance. If it is terminal
use the size floor. Thus any child meeting a small ball of
radius \(c'\ell_S\) about \(x\) has area \(\gtrsim\ell_S^2\).
Every child meeting \(B(x,C\ell_S)\) has diameter
\(O(\ell_S)\), by the upper side comparison and the Lipschitz
bound on distance. The support bound follows from
\eqref{eq:ray-support-self-contained}. On near children,
\(\int_{B(x,C\ell_S)}|x-y|^{-1}dy=O(\ell_S)\)
divided by the child area gives the desired supremum.
Farther away use the pointwise bound. For the difference, put
\(d=|x-x'|\). When \(d\) is a small fraction of \(\ell_S\),
integrating the two singularities in the \(2d\) balls and the
derivative elsewhere gives \(O(d(1+\log(\ell_S/d)))\) over
the near region. Divide by the near child area. For far
children the derivative estimate gives \(O(d/\ell_S^2)\)
pointwise. Use \(v(1+\log(1/v))\lesssim_\beta v^\beta\).
When \(d\) is comparable to \(\ell_S\), the single-point
estimates suffice. The density bound follows from the
neighborhood comparability. The resulting \(L^2\) kernel
bounds are uniform. Subtracting \(P_Q\) contracts \(L^2\).

Here is the Gaussian envelope argument, also used for the
fine split below. For any exponent \(\beta\) between \(\alpha\)
and one the \(L^2\) bounds control the canonical increment
distance by \(C(|x-x'|/\ell_S)^\beta\). In the rescaled
neighborhood, take nets of mesh \(2^{-j}\) and cardinality
at most \(C2^{2j}\). Compare points in two successive nets
within distance \(C2^{-j}\), and also compare neighboring
points at each net level. The expected maximum of the
absolute Gaussian differences of these pairs is at most
\(C2^{-j\beta}\sqrt{1+j}\), by the Gaussian tail and a union
bound. Telescoping from the level of the separation of two
points bounds the expected scaled \(\alpha\)-H\"older
seminorm by the convergent sum
\(C\sum_{j\ge0}2^{-j(\beta-\alpha)}\sqrt{1+j}\).
It also constructs a continuous version. The expected
supremum is bounded by the same sum and the expected
absolute value at one point. The density tests are constant
on only a bounded number of cells in the neighborhood.
Finally the three envelope terms are suprema of linear
Gaussian tests, with signs or unit vector directions
included. Their test norms are uniformly bounded in
\(L^2\), so the sum is Lipschitz in the Gaussian coefficients
with bounded constant. Gaussian concentration proves the
tail, and the definition gives evenness.
\end{proof}

\begin{lemma}[Sharp covariance for a uniform fine split]
\label{lem:sharp-fine-kernel}
Let \(Q\) be a square subdivided uniformly into \(R\times R\)
squares of side \(\ell\). There is a zero-normal inverse
for the fresh contrast with chart covariance at most
\[
 \left(\sqrt{\frac{\log R}{4\pi}}+C\right)^2 I .
\]
For \(x,x'\) in a fine child or its small \(\ell\)-neighborhood
inside \(Q\), the canonical \(L^2\) distance of the two
evaluations is at most \(C_\beta(|x-x'|/\ell)^\beta\) for
every fixed \(\beta<1\). An even envelope controlling the
supremum and scaled \(\alpha\)-H\"older seminorm of the current,
divided by \(\sqrt{1+\log R}\), has a uniform Gaussian tail.
\end{lemma}

\begin{proof}
Choose a smooth radial probability density \(\zeta\) supported
in the annulus \(1/2<|z|<1\). For \(y\in Q\), put
\(a_y=c\,d_Q(y)\) for a small fixed \(c>0\) and set
\(\zeta_y(z)=a_y^{-2}\zeta((z-y)/a_y)\). Define
\[
\begin{aligned}
 R_y(x)&=\frac{x-y}{2\pi |x-y|^2}
       \left(1-\int_{B(y,|x-y|)}\zeta_y(z)\,dz\right),\\
 B(x,y)&=\int_Q K_Q(x,z)\zeta_y(z)\,dz,\qquad
 \widetilde K_Q(x,y)=R_y(x)+B(x,y).
\end{aligned}
\]
Radial integration shows that \(R_y\) has divergence
\(\delta_y-\zeta_y\). Its support is in \(B(y,a_y)\).
The divergence of \(B\) is \(\zeta_y-\eta_Q\), by the ray
identity. Thus \(\widetilde K_Q\) has divergence
\(\delta_y-\eta_Q\) and zero normal flux. Both summands vanish
for \(|x-y|>C d_Q(x)\). For \(B\) this follows from the ray
support bound and \(a_y\le c(d_Q(x)+|x-y|)\). The same
inequality and the radial support give the assertion for \(R_y\).

On this support \(B\) and \(\nabla_x B\) are bounded by
\(C/d_Q(x)\) and \(C/d_Q(x)^2\), respectively. When
\(|x-y|\ge d_Q(x)/4\), the smallness of \(c\) keeps the
support of \(\zeta_y\) a distance comparable to \(|x-y|\)
from \(x\), so the estimates follow from the ray bounds.
Otherwise \(d_Q(y)\asymp d_Q(x)\), and integration of
\(|x-z|^{-1}\) against \(\zeta_y\) gives the first bound.
For the derivative, replace \(\nabla_x K_Q\) by
\((\nabla_x+\nabla_z)K_Q-\nabla_z K_Q\) and integrate the
second term by parts against \(\zeta_y\). The simultaneous
derivative is bounded by \(C/(D|x-z|)\), and the other term
uses \(\|\nabla\zeta_y\|_\infty\lesssim d_Q(x)^{-3}\).
This proves the second bound. The radial term and its
derivative are bounded by \(C/|x-y|\) and \(C/|x-y|^2\).
It follows that the full kernel satisfies the same pointwise
bounds. Its \(L^1\) norm in the source variable is
\(O(d_Q(x))\), by support and the singular bound, so currents
of bounded sources extend continuously by zero on \(\partial Q\).

Let \(P\) now average on the uniform fine grid and fix a unit
direction \(e\). On \(|x-y|\le C_1\ell\), the projected
singular kernel has \(L^2\) norm \(O(1)\). Indeed only \(O(1)\)
fine cells meet that disk, their area is \(\ell^2\), and the
\(L^1\) norm in the disk is \(O(\ell)\). The entire regular
part \(B(x,\cdot)\) has \(L^2\) norm \(O(1)\) by its support
and pointwise bound. On
\(C_1\ell<|x-y|<c_1d_Q(x)\), with \(c_1\) sufficiently small,
the radial kernel is the free-space vector
\((x-y)/(2\pi|x-y|^2)\). This is because \(d_Q(y)\asymp d_Q(x)\)
there and \(\zeta_y\) vanishes on \(B(y,a_y/2)\).
The whole annulus lies inside \(Q\). Its squared \(L^2\) norm
in direction \(e\) is
\[
 \int_{C_1\ell}^{c_1d_Q(x)}
       \int_0^{2\pi}\frac{\cos^2\vartheta}{4\pi^2 r^2}
              r\,d\vartheta\,dr
 = \frac{1}{4\pi}\log\!\left(\frac{c_1d_Q(x)}{C_1\ell}\right).
\]
The remaining radial support lies in an annulus with bounded
ratio of outer and inner radii, giving \(L^2\) norm \(O(1)\).
If the displayed annulus is empty the bounds on the near
part and remaining support still give \(O(1)\).
Orthogonal projection contracts \(L^2\), and
\(d_Q(x)/\ell\le CR\). The triangle inequality for these
pieces therefore gives the stated covariance for
\(W((P-P_Q)\widetilde K_Q(x,\cdot))\).

For the difference let \(d=|x-x'|\le c'\ell\). On a disk
of radius \(C_1\ell\) containing the two poles, the two-pole
\(L^1\) bound is \(O(d(1+\log(\ell/d)))\). Averaging on the
\(O(1)\) intersecting cells gives \(L^2\) norm bounded by this
quantity divided by \(\ell\). Outside the disk the derivative
bound gives squared \(L^2\) norm \(O(d^2/\ell^2)\) before
projection. For \(d\) comparable to \(\ell\) the same split
uses the separate \(L^1\) singular bounds on the near region
and the derivative bound on its complement. This proves the
canonical increment estimate. The net argument of
Lemma~\ref{lem:coarse-envelope} now gives an expected oscillation
envelope \(O(1)\) before division, and the expected supremum is
\(O(\sqrt{1+\log R})\) by adding one evaluation. The normalized
linear tests have uniformly bounded \(L^2\) norm, so Gaussian
concentration and evenness hold as claimed.
\end{proof}

\section{Local flow estimates and repair of omitted sources}
\label{sec:sc-repairs}

\subsection{Uniform capacity on a square mesh}

Let \(D\) be a square made of \(K^2\) grid cells, with \(K\) a
power of two. Orient the internal grid edges and let \(d q\)
denote the outgoing load minus the incoming load at each cell.
Put
\[
 \mathcal B_D(g)=\min_{dq=g}\max_e|q_e|,\qquad \sum_vg_v=0 .
\]
There are no loads on the exterior boundary. For a cell set \(S\)
write \(\partial S\) for the internal edges between \(S\) and its
complement. Then
\begin{equation}
 \mathcal B_D(g)=
  \max_{\varnothing\ne S\ne D}\frac{|\sum_{v\in S}g_v|}{|\partial S|}.
 \label{eq:sc-box-cut}
\end{equation}
The upper bound in \eqref{eq:sc-box-cut} follows by separation of \(g\)
from the divergence of the edge cube of the proposed capacity.
Indeed, the support function of this image at a vector \(u\) is
the capacity times \(\sum_{v\sim w}|u_v-u_w|\). Subtract \(\min_v u_v\)
from \(u\) and integrate its superlevel sets. The integral of their
cut sizes is \(\sum_{v\sim w}|u_v-u_w|\), and the integral of their
signed \(g\)-masses is \(\sum_vu_vg_v\). Thus the cut inequalities
give every support-function inequality for the image. The converse
is obtained by summing \(dq=g\) over \(S\).

For \(|S|\le K^2/2\), if no row or column is full, the product of
the numbers of occupied rows and columns is at least \(|S|\).
Each occupied row or column contains a cut edge in its own
direction. If a row is full, every column meets \(S\) and at
least \(K/2\) columns are mixed, and the argument for a full
column is the same with axes exchanged. Hence
\[
 |\partial S|\ge c\sqrt{\min(|S|,|D\setminus S|)} .
\]
It follows from Cauchy--Schwarz and \eqref{eq:sc-box-cut} that
\(\mathcal B_D(g)\le C\|g\|_2\), and the triangle inequality gives
the same Lipschitz bound for \(\mathcal B_D\).

When the cell width is \(w\), linearly interpolate each component
between its two normal faces inside each cell, leaving it
constant in the other coordinate. The resulting current has
component suprema at most the maximum face load, divergence
\(g_v/w\) on cell \(v\), and zero normal trace on the box boundary.
The normal traces agree across internal faces. A separate
coordinate primitive will also be useful. In a rectangle
\(I\times J\) with a bounded mean-zero density \(f\), put
\(r(y)=\int_I f(x,y)\,dx\) and use the two components
\begin{equation}
 J_x(x,y)=\int_{\inf I}^x\left(f(t,y)-\frac{r(y)}{|I|}\right)dt,
 \qquad
 J_y(x,y)=\frac{1}{|I|}\int_{\inf J}^y r(t)\,dt .
 \label{eq:sc-row-primitive}
\end{equation}
These have divergence \(f\) and zero outer normal trace. In
particular, on a square of side \(w\) both bounds are
\(2w\|f\|_\infty\).

\begin{lemma}\label{lem:sc-box-gauss}
Let \(G\) be a centered Gaussian mean-zero cell vector in \(D\)
with covariance at most the orthogonal projection \(P\) onto
mean-zero vectors. Then
\[
 \E\mathcal B_D(G)\le C(1+\log K),\qquad
 \PP\{\mathcal B_D(G)>C(1+\log K)+u\}\le e^{-c u^2}.
\]
\end{lemma}
\begin{proof}
For independent standard Gaussian cell coefficients \(Z\), telescope
\(PZ\) into dyadic child-average contrasts. In each parent of
side \(l\) grid cells the four child values have standard
deviations at most \(C/l\). Route their contrast by row partial
sums after removing the row means, then use a vertical partial
sum of the row means. The parent is balanced, so both outer
normal loads vanish. Each face load is linear in the four
child values, with coefficients bounded by \(Cl\), and
therefore has variance at most \(C\). The levels of the dyadic
contrast are independent by orthogonality. Summing the linear
flows gives face variance at most \(C(1+\log K)\), and
the Gaussian tail bound on the \(O(K^2)\) faces gives expected
maximum at most \(C(1+\log K)\). For covariance dominated by
\(P\), add an independent centered Gaussian of the covariance
difference and apply convexity to get the same mean bound.
The Euclidean Lipschitz bound from \eqref{eq:sc-box-cut}
and Gaussian concentration give the tail.
\end{proof}

We use Gaussian concentration in its Lipschitz form
\[
 \E e^{t(F(Z)-\E F(Z))}\le e^{t^2 L^2/2}
\]
for an \(L\)-Lipschitz function \(F\) of a standard Gaussian vector.
One derivation is to set \(P_tg(x)=\E g(e^{-t}x+
\sqrt{1-e^{-2t}}Z)\). Gaussian integration by parts gives
\(-\partial_t\E[(P_tg)\log(P_tg)]=\E[|\nabla P_tg|^2/P_tg]\).
Commuting the gradient through \(P_t\) and applying Cauchy--Schwarz
bounds the right side by \(e^{-2t}\E(|\nabla g|^2/g)\).
Integrating in \(t\) gives the Gaussian entropy bound with
constant \(1/2\). With \(g=e^{sF}\) it bounds
\(\partial_s(s^{-1}\log\E e^{sF})\) by \(L^2/2\), which
gives the displayed moment bound upon integration. Approximation
extends the calculation from smooth bounded functions to Lipschitz
\(F\).

\begin{lemma}\label{lem:sc-box-emp}
Let \(g=\sum_i \sigma_i(e_{V_i}-p_i)\), where the \(V_i\)
are independent cell labels, \(p_i=\E e_{V_i}\), and
\(\sigma_i\in\{-1,1\}\) are fixed. If the total mean occupancy
per cell is at most \(C\mu\), \(\mu\asymp m\), and \(K\le m^C\),
then
\[
 \E\mathcal B_D(g)\le C'\sqrt\mu(1+\log K),\qquad
 \PP\{\mathcal B_D(g)>C_A\mu\}\le e^{-A m}
\]
for each fixed \(A\), with the constant chosen accordingly.
These estimates also apply conditionally with the same parameters.
\end{lemma}
\begin{proof}
Use an independent replica \(g'\) and Jensen to bound the
mean by that of \(g-g'\). The paired label differences
\(\sigma_i(e_{V_i}-e_{V_i'})\) can be symmetrized by independent
signs, since the ordered pairs are exchangeable. Given the two
label arrays, replace the signs by standard normals at a
constant dilation \(1/\E|Z_1|\). This follows by conditioning
on the Gaussian signs and applying Jensen over their absolute
values. If \(D_*\) is the largest pooled occupancy of a cell
in the two arrays, the conditional covariance of this Gaussian
sum is at most \(C D_*P\), by
\(\sum_i (u_{V_i}-u_{V_i'})^2\le 2D_*\|Pu\|_2^2\).
Each pooled cell count \(N_v\) has
\(\E e^{N_v}\le\exp(C'\mu)\) by the product moment of its
independent indicators. Thus
\(\E D_*\le C'\mu+\log(K^2)\), and
\(\PP(D_*>C_A'\mu)\le e^{-2A m}\) after enlarging the constant.
The conditional Gaussian mean bound and Jensen give the
claimed unconditional mean. On \(D_*\le C_A'\mu\), apply
the sign and covariance comparisons to the convex function
\(\exp(t\mathcal B_D)\). The Gaussian moment bound gives
\[
 \exp\!\left(C_A''t\sqrt\mu(1+\log K)+C_A''t^2\mu\right).
\]
A fixed positive tilt and the pooled occupancy estimate give
\(\PP(\mathcal B_D(g-g')>C_A'''\mu)\le \tfrac12 e^{-A m}\)
for a suitable constant. Finally for \(u=C_A'''\mu\)
the probability that \(\mathcal B_D(g')\le u\) is at least
\(1/2\) by the mean bound. Independence and the triangle inequality
then give the asserted bound for \(\mathcal B_D(g)>2u\).
\end{proof}

Here and below we use the following elementary count estimates.
For a sum \(S\) of independent indicators of total mean \(\mu\),
its moment generating function is bounded by
\(\E e^{t(S-\mu)}\le\exp(\mu(e^t-1-t))\).
This is exact for a Poisson variable of mean \(\mu\). For
\(u\ge0\), exponential tilts of both signs give
\(\PP(|S-\mu|>u)\le 2e^{-c\min(u^2/\mu,u)}\).
For example, use \(e^t-1-t\le C t^2\) for \(|t|\le1\)
and choose \(|t|\) a sufficiently small constant times
\(\min(1,u/\mu)\). The same argument applied to independent
centered random variables \(X_i\) with \(|X_i|\le B\) gives
\(\PP(|\sum_i X_i|>u)\le 2e^{-c\min(u^2/v,u/B)}\)
for \(v=\sum_i\E X_i^2\), by using
\(e^{tX_i}-1-tX_i\le C t^2 X_i^2\) for \(|t|\le 1/B\).
In particular the estimates apply to each cell in a
multinomial count, and conditionally to any independent
labels left by conditioning on the coarser counts.

\subsection{Joint tails for coarse envelopes}

At a predictable generation let \(S\) range over the Whitney
children of the old parents. Write \(l_S\) for intrinsic
child size and \(E_S\) for the coarse envelope in the parent
chart. We use the local kernel and packing bounds of
Lemmas~\ref{lem:local-geometry} and~\ref{lem:coarse-envelope}.
In particular the envelope is a sum of a fixed number of
suprema of absolute normalized Gaussian tests, each with
a kernel of size at most \(C/l_S\) and supported within
\(Cl_S\) of \(S\), before the projection away from parent
constants. The single-child mean is bounded uniformly.

\begin{lemma}\label{lem:sc-joint}
Conditionally on the old history, for each predictable finite
subset \(\mathcal A\) of attempted children and sufficiently
large \(v\),
\[
 \PP\{E_S>v\ \text{for all } S\in\mathcal A
                  \mid\text{old history}\}
 \le \exp(-c|\mathcal A|v^2).
\]
\end{lemma}
\begin{proof}
Consider one normalized test kernel \(k_S\) selected for
each \(S\in\mathcal A\). For overlapping supports in
the same chart with \(l\le l'\) their scalar product is
at most \(Cl/l'\) in absolute value, by the support-area
bound. For a fixed child of size \(l'\), the number of
overlapping kernels with size \(l<l'\) is \(O(l'/l)\).
Indeed their disjoint child areas lie in a boundary
strip of width \(Cl\) within a ball of radius \(Cl'\).
For each \(l'\ge l\), at most \(O(1)\) of the larger
supports can overlap a fixed size-\(l\) support, by disjoint
areas comparable to \(l'^2\). Weighted Schur with weights
\(l_S^\eta\), \(0<\eta<1\), bounds the Gram matrix uniformly.
In the sum from smaller children the factor is
\((l/l')^\eta\), and in the sum from larger children it is
\((l/l')^{1-\eta}\). The resulting geometric sums are finite.
The same argument for comparable scales is just area packing.
Thus for arbitrary coefficients \(a_S\),
\(\|\sum_S a_S k_S\|_2^2\le C\sum_S|a_S|^2\).
Within each parent projection away from its constant
is a common orthogonal projection, which decreases this
norm. The projected tests from distinct parents have
disjoint supports. Taking suprema over the choices of
one signed test per child and summing a bounded number
of these suprema, \(\sum_{\mathcal A}E_S\) is consequently
\(C\sqrt{|\mathcal A|}\)-Lipschitz. Its conditional mean is
\(O(|\mathcal A|)\). Gaussian concentration at
\(|\mathcal A|v\) proves the result.
\end{proof}

\subsection{Coarse clipping and exact sparse repair}

Fix a threshold \(E_0\). Take a Lipschitz cutoff \(\chi_S\)
that is one on \(S\) and is supported in the small chart
neighborhood defining \(E_S\), with
\(|\nabla\chi_S|\le C/l_S\) in that chart. For example, use
the positive part of \(1-\operatorname{dist}(\cdot,S)/(c l_S)\)
with \(c\) sufficiently small. Call \(S\) exceptional if
\(E_S>E_0\) and multiply its parent's whole current by
\[
  1-\chi_*=\prod_{S\ {\rm exceptional}}(1-\chi_S).
\]
This multiplier is even in the new noise and is between zero
and one. Bounded comparable overlap implies that on each
nonexceptional child the supremum and scaled H\"older bounds
for the retained current are \(CE_S\), and on exceptional
children the retained current is zero. To verify the H\"older
bound within a child, use the product difference inequality
for factors in \([0,1]\), the uniformly bounded number of
cutoffs meeting the child, and their comparable
\(O(1/l_S)\) gradient bounds. These are the bounds on the child used in the evolution.
At any fixed point the even multiplier only decreases the conditional
linear exponential moment, by averaging over the signs
of the entire new Gaussian vector and using that
\(\cosh(v)\) increases with \(|v|\).

Order the exceptional children within a parent and
decompose \(\chi_*V=\sum_S\varphi_S V\), with
\(\varphi_S=\chi_S\prod_{T<S}(1-\chi_T)\).
For the current multiplied by a deterministic scale \(\sigma\),
the omitted source \(f_S=\sigma\Div(\varphi_S V)\) is bounded
by \(C\sigma E_S/l_S\) in its chart neighborhood. The
product rule uses the density test in \(E_S\) and
\(|\nabla\varphi_S|\le C/l_S\), again by comparable overlap.
The parent current has zero normal trace, so the product
rule across its boundary introduces no measure on that
boundary. Thus each source is a genuine bounded
function of integral zero, supported in the physical
image of a chart ball of radius \(Cl_S\) intersected
with the parent.

\begin{lemma}\label{lem:sc-sparse}
Suppose \(m=\log(1/h)\), there are at most \(Cm\) generations,
all physical charts have determinant one and distortion bounded
by a fixed power of \(m\), and \(h\le l_S\le m^C\).
Assume the domain is an area-one rectangle with both sides and
their reciprocals at most \(m^C\). Take
\(\log m=o(E_0^2)\), \(E_0^2\le m\), and set
\[
 D_0=m/E_0^2,\qquad T_0=\sigma\sqrt m .
\]
For any \(1\le B\le m^C\), with probability \(1-o(1)\)
the sum of all omitted sources has an exact current with
zero outer normal trace, supported on a set of area at
most \(\exp(-cE_0^2)\), and satisfying
\[
 \|J_y\|_\infty\le C T_0 D_0/B,\qquad
 \|J_x\|_{\psi_2}\le C T_0 D_0 B/E_0 .
\]
\end{lemma}
\begin{proof}
A conditional single-child tail and a union bound over
the at most \(Cm h^{-2}\) attempted children up to constants
give \(E_S\le C\sqrt m\) everywhere with probability
\(1-o(1)\). Thus the source in its support ellipse
satisfies
\[
 \|f_S\|_\infty\le C T_0/l_S,\qquad
 \|f_S\|_1\le C T_0 l_S .
\]
If \(A\) is the chart map and \(G=AA^{\mathsf T}\), the
horizontal and vertical maximal chord lengths of the ellipse
are bounded by \(Cl_S/\sqrt{G_{yy}}\) and
\(Cl_S/\sqrt{G_{xx}}\). This follows by imposing zero
orthogonal physical coordinate on a chord in the equation
\((z-z_0)^{\mathsf T}G^{-1}(z-z_0)\le C l_S^2\) and using
\(\det G=1\). Consequently the absolute horizontal and
vertical line masses \(A_x,A_y\) of \(f_S\) satisfy
\(A_xA_y\le C T_0^2\).

Around each candidate ellipse choose a predictable axis box
inside the domain with both sides at least \(B l_S\).
For instance enclose the ellipse portion in the domain
in its axis bounds, enlarge each side by \(B l_S\)
as necessary, and slide the intervals inside the domain.
The side requirements are at most \(m^{C'}l_S\) by
the chart cap. If the box cannot fit, declare failure
only if \(S\) is exceptional. Such a candidate has size
at least \(m^{-C'}\), so there are only polynomially
many per generation by area packing. The probability
of a failure is \(o(1)\) by the conditional envelope tail.

When \(A_x\le T_0 B\), apply \eqref{eq:sc-row-primitive}
in the candidate box. Its horizontal supremum is at
most \(2A_x\) and its vertical one at most
\(\|f_S\|_1/|I|\le CT_0/B\). If \(A_x>T_0 B\),
then \(A_y\le CT_0/B\), and applying the column version
gives the same bounds with a changed constant, since
the horizontal supremum of the spread mass primitive
is at most \(CT_0/B\le CT_0 B\). Extend the box
currents by zero. Their zero normal traces preserve
the divergence.

To bound overlap, fix a point and pad the candidate
boxes by \(O(h^2)\) on all sides. The number \(n\) of
these boxes containing the point in a generation is at
most \(m^{C''}\). In each intrinsic dyadic size class their
disjoint child areas are comparable to the squared size
and lie within polynomial times that size of the point.
The padding preserves this bound as \(l_S\ge h\).
If \(I_S\) is the exceptional indicator, expand
\[
 \E\exp\!\left(\theta\sum_S I_S\right)
 =\E\prod_S(1+(e^\theta-1)I_S)
 \le \exp(ne^{-cE_0^2/2}),\qquad
 \theta=cE_0^2/2,
\]
conditionally on the old history, by Lemma~\ref{lem:sc-joint}.
Iterate through the generations. The probability of overlap
larger than \(C'D_0\) at the point is at most
\(\exp(-c'C'm+o(1))\). A union bound over a mesh of
spacing \(h^2\) and the padding give a uniform cap
\(O(D_0)\) for a sufficiently large \(C'\).
The sum of all candidate box areas in a generation
is bounded by a fixed power of \(m\) by disjoint
child areas. The expected total area of exceptional
boxes is thus \(m^{C'''}e^{-cE_0^2}\), by the single
envelope tail. Markov's inequality gives the stated
area bound with high probability. A function bounded by
\(M\) on a set of area \(p\) and zero off it has Luxemburg
norm at most \(M/\sqrt{\log(1+1/p)}\). Apply this with
\(M=CT_0 D_0 B\) and the area bound. The overlap cap
also gives the vertical supremum bound.
\end{proof}

\section{The sharp Gaussian lower bound}
\label{sec:sharp-constant-lower}

Set \(m=\log(1/h)\), and let \(R\) be a dyadic integer with
\(\log R\sim(\log m)^2\). We construct a random Lipschitz test
potential \(u\) by adding local Poisson solutions on successively
smaller squares. Its vertical derivative is positive with integral
one. Changing spatial measure by this derivative makes the ratio of
horizontal to vertical derivatives a martingale. We control that
ratio by stopping stages of a state-dependent variance clock.

\subsection{One Poisson refinement}

On \(S=(-1/2,1/2)^2\), subdivide the lattice squares of side
\(R^{-1}\) into four triangles by joining their corners to their
centers. Let \(I_R\) be nodal affine interpolation. Take a
dihedrally invariant smooth cutoff \(\eta\) which is zero on a
boundary strip of width \(m^{-A_0}/2\) and one outside the strip of
width \(m^{-A_0}\). Its derivatives of order \(k\) are bounded by
\(C_k m^{kA_0}\). Let \(\phi_k\) run over an orthonormal real
Fourier basis in the band
\[
 m^{A_1}\le |k|\le Rm^{-A_1},\qquad \lambda_k=4\pi^2|k|^2 ,
\]
where both sine and cosine are included for each frequency pair,
and \(A_1\) is a sufficiently large fixed constant in terms of
\(A_0\). Denote these cutoffs by \(K_{\min}\) and \(K_{\max}\).
Put \(g_k=(I-\Pi_{\rm aff})I_R\phi_k\), where \(\Pi_{\rm aff}\)
is orthogonal projection onto affine functions. Define \(f_k\)
by subtracting from \(I_R(\eta\phi_k)\) a linear combination of
\(I_R(\eta_0p)\), \(p=1,x_1,x_2\), so that \(f_k\) is orthogonal to
all affine functions. Here \(\eta_0\) is a fixed dihedrally invariant
nonnegative smooth cutoff supported strictly inside \(S\) and
positive on an open subset. The \(3\times3\) moment matrix
converges to the positive definite Gram matrix with weight
\(\eta_0\). This defines the correction uniquely for large \(m\),
and it commutes with the dihedral action. Both \(g_k\) and \(f_k\)
are piecewise affine, and \(f_k\) vanishes near \(\partial S\).

For white noise \(\mathcal W\) on \(S\), let
\[
 F=\sum_k \lambda_k^{-1}\mathcal W(g_k)f_k,\qquad
 V_R=\frac{\log R}{4\pi}.
\]
The coefficient covariance matrix is \(I+o(1)\) in operator norm.
Indeed, on each shape-regular triangle the interpolation error in
\(L^2\) is at most \(CR^{-2}\) times the \(L^2\) norm of the
second derivatives. To obtain the unit-scale estimate, subtract
the affine function making the mean and mean first derivatives
zero, apply Poincar\'e twice, and use the embedding of \(H^2\)
in \(C^0\) in two dimensions to control the nodal values. Summing
over triangles gives an error operator of norm
\(O((K_{\max}/R)^2)\) on the Fourier band. The affine projection
of the Fourier band has norm \(O(K_{\min}^{-1/2})\) because the Fourier
coefficients of each coordinate lie on its coordinate axis
and are \(O(|k|^{-1})\), whose squares sum to
\(O(K_{\min}^{-1})\). These two estimates give the claimed Gram
matrix bound.

Away from triangle boundaries, the weighted coefficient vector
of \(\nabla F\) differs in \(\ell^2\) by \(o(1)\) from the vector
\((\lambda_k^{-1}\eta\nabla\phi_k)_k\). The interpolation
gradient error in a single frequency is bounded by \(CR^{-1}\)
times the supremum of its second derivatives, by Taylor's
formula at the triangle vertices. Here that gives
\[
 |\nabla I_R(\eta\phi_k)-\nabla(\eta\phi_k)|
 \le CR^{-1}(|k|^2+m^{A_0}|k|+m^{2A_0}).
\]
The weighted \(\ell^2\) error is \(O(K_{\max}/R)\) as
\(K_{\min}\gg m^{2A_0}\). The term containing \(\nabla\eta\)
costs \(O(m^{A_0}/K_{\min})\) in the same norm.
The moment corrections have uniformly bounded coefficients
since the uncorrected output functions are uniformly bounded
and the moment matrix has bounded inverse. Their weighted
\(\ell^2\) gradient contribution is \(O(K_{\min}^{-1})\).
The same estimates imply \(\|\nabla f_k\|_2\le C|k|\) and
\(\|f_k\|_2\le C\). In the ideal periodic gradient, the
sine-cosine pair sums and lattice symmetry give
\[
 \sum_k\lambda_k^{-2}\nabla\phi_k(x)\nabla\phi_k(x)^{\mathsf T}
  =\left(\frac{1}{4\pi}\log(K_{\max}/K_{\min})+O(1)\right) I .
\]
In fact the trace equals
\((4\pi^2)^{-1}\sum_{K_{\min}\le |k|\le K_{\max}}|k|^{-2}\)
with the sum over lattice vectors. Comparing with the radial
integral gives \(2\pi\log(K_{\max}/K_{\min})+O(1)\) for the
lattice sum, since the derivative of \(|x|^{-2}\) is integrable
away from the origin. Thus the covariance of \(\nabla F\) is
bounded by \(CV_R I\) and is \((1+o(1))V_R I\) outside the
boundary strip. The integral in \(x\) of the covariance is a
scalar matrix by dihedral symmetry.

The expected unprojected pairing is
\[
 \E\sum_k\lambda_k^{-1}\mathcal W(g_k)\mathcal W(f_k)
  =\sum_k\lambda_k^{-1}\langle g_k,f_k\rangle=(2+o(1))V_R .
\]
For the last equality,
\(\langle g_k,f_k\rangle=\langle\phi_k,\eta\phi_k\rangle+o(1)\)
uniformly. For the interpolation errors use \(K_{\max}/R=o(1)\)
and the preceding bounds. For the moment correction, each
\(I_R(\eta_0p)\) has inner product \(o(1)\) with each band
mode by integration by parts and its \(O(R^{-2})\)
interpolation error. Summing the sine-cosine pairs replaces
\(\langle\phi_k,\eta\phi_k\rangle\) by \(\int_S\eta=1+o(1)\).
The same lattice sum proves the formula.

Clip \(F\) to zero if \(\|\nabla F\|_\infty>C\log R\). This
event has probability \(O(R^{-10})\) with a sufficiently large
fixed \(C\), by a Gaussian tail and a union bound over the
\(4R^2\) triangles. Let \(\widehat F\) be the clipped function.
Its integrated covariance is still a scalar matrix \(v I\),
since the clipping event is invariant under the dihedral
group. Gaussian fourth moments and Cauchy--Schwarz show
\(v=(1+o(1))V_R\). Set \(U=\widehat F/\sqrt v\).
For an independent uniform point of \(S\), the gradient
\(\xi=\nabla U\) is centered with covariance \(I\) and satisfies
\begin{equation}\label{eq:lower-innovation}
 |\xi|\le C\sqrt{\log R},\qquad
 \int_S\nabla U=0\quad\text{for every realization}.
\end{equation}
The clipped pairing has expectation
\((1+o(1))\sqrt{\log R/\pi}\). Indeed, the unconditional \(L^2\)
norm of the pairing before division by \(\sqrt v\) is
\(O(\log R)\), by Gaussian fourth moments of each factor,
Minkowski, and \(\sum_k\lambda_k^{-1}=O(\log R)\).
Cauchy--Schwarz bounds the clipping error by \(o(\log R)\).

\subsection{The tilted stage calculation}

Fix constants \(\gamma,B>0\) and \(0<\varepsilon<1\).
Put \(\kappa=1/20\) and choose
\[
 J=\left\lfloor\frac{\kappa\log m}{8\log\log m}\right\rfloor,\quad
 q=\frac{\gamma}{J},\quad L=\log(1/q),\quad
 S_j=m^\kappa L^{4j},\quad
 b_j=\frac{B}{\sqrt{S_j}},\quad \ell_j=qb_j,\quad A_j=\frac{S_j}{2L}.
\]
Then \(m^\kappa\le S_j\le m^{2\kappa}\) for large \(m\), and
\(L\sim\log\log m\). Let \(n\sim m/\log R\) be the maximum
number of refinement steps, fixed below. In stage \(j\) take
\begin{equation}\label{eq:lower-stage-clock}
 a_j(x)=\frac{A_j}{x^2+\ell_j^2},\qquad
 r_j=\frac1{n\lambda_j},\qquad
 du=r_j a_j(x)^{-1/2},\qquad ds=r_j a_j(x)^{1/2},
\end{equation}
where \(\lambda_j\asymp_{\gamma,B,\varepsilon}L^{-1/2}\) is
specified below. These are the variance clock increments and
\(\sqrt{du\,ds}=r_j\) is the resource paid at each step.

Let \((X,Y)\) denote the gradient of the current potential on
a branch, initially \((0,1)\), and write \(M=X/Y\).
In stage \(j\), the coordinate \(x\) in \eqref{eq:lower-stage-clock} is \(M\) minus its
value at the entrance to the stage. On an active step use
\[
 \binom{\Delta X}{\Delta Y}
  =Y\sqrt{r_j}
   \begin{pmatrix}a_j^{-1/4}&M a_j^{1/4}\\0&a_j^{1/4}\end{pmatrix}\xi .
\]
Writing \(\beta=\Delta Y/Y\) and
\(c=(\Delta X-M\Delta Y)/Y\), the ordinary centered innovations
have variances \(ds\) and \(du\), respectively. Their absolute
values are bounded by \(C\sqrt{ds\log R}\) and
\(C\sqrt{du\log R}\). Each updated block preserves the integral
of \(Y\), by \eqref{eq:lower-innovation}. The pathwise spatial law tilted by the final
\(Y\) is a probability measure because the initial integral is
one. We write \(\nu\) for the joint law obtained by applying this
tilt to the uniform spatial branch for each noise realization.

On a retained block, the conditional change of measure is
\(1+\beta\). To see this even with geometry chosen adaptively,
condition on the earlier reveals and the old branch cell. The
integral of the final \(Y\) over any current cell equals its
current constant \(Y\) times its area, since every later
perturbation lies inside a subcell and has zero integrated
gradient. If a retiling discards an area fraction \(\theta\)
of the cell, that region keeps the old \(Y\), and the retained
squares each have spatial average \(Y\). The discard probability
is thus exactly \(\theta\) under either law, conditional on the
old data. Conditional on retention, the branch is uniform in
its new square before the next innovation. For each fixed fresh
field, its total future spatial weight in that square is the
old \(Y\) times the area. The marginal fresh Gaussian field
therefore retains its original law, and the child selection
density relative to uniform is \(1+\beta\). Consequently
\[
 \Delta M=\frac{c}{1+\beta},\qquad
 \E_\nu[\Delta M\mid\mathrm{past}]=0,\qquad
 \E_\nu[(\Delta M)^2\mid\mathrm{past}]=(1+o(1))du
\]
on retained steps. In fact the last conditional expectation
equals \(\E[c^2/(1+\beta)\mid\mathrm{past}]\) under the original law.
Before \(|x|\) exits \(b_j\), \eqref{eq:lower-stage-clock} gives
\[
 ds\le C S_j/(q n),\qquad
 du/b_j^2\le C L/n,\qquad
 |\beta|\le C\sqrt{S_j\log R/(q n)}=O(m^{-c_0})
\]
for an absolute \(c_0>0\). The ratio jumps are bounded by
\(D_j=C b_j\sqrt{L\log R/n}\). Uniformly in \(j\),
\(D_j/\ell_j=O(m^{-c_0})\) and
\(\sum_j D_j=o(b_J/\sqrt{\log J})\), by decreasing \(c_0\)
if necessary. These bounds also hold for synthetic continuations
using the same law of \(\xi\), tilted by \(1+\beta\) in each step.

For the moment consider an untruncated stage, with no time,
geometry or clock cap. Its relative ratio \(x_k\) starts at
zero, and the discrete local time is
\[
 L_k^0=|x_k|-\sum_{i=1}^k
         \sgn_+(x_{i-1})(x_i-x_{i-1}),\qquad \sgn_+(0)=1 .
\]
Its increments are nonnegative and at most \(2D_j\), by the
convexity and Lipschitz bounds for \(|x|\). A positive increment
ends within \(D_j\) of zero. Stop the stage on \(L_k^0\ge\ell_j\)
(a success) or \(|x_k|\ge b_j\) (a failure \(F_j\)). These
conditions cannot occur in the same step for large \(m\).
The stopped position is bounded by \(b_j+D_j\), and the
conditional variance before stopping is bounded below by a
positive deterministic number for fixed \(m,j\). Applying the
stopped square identity at bounded times shows that the stopping
depth has finite mean. Optional stopping for the Tanaka identity
gives \(\E_\nu|x_T|=\E_\nu L_T^0\). Thus the conditional failure
probability \(p_j\) at entry satisfies
\[
 p_j b_j\le\ell_j+2D_j,\qquad
 (1-p_j)\ell_j\le p_j(b_j+D_j)+(1-p_j)D_j .
\]
In particular \(p_j=(1+o(1))q\) uniformly in the entry history.
From an intermediate state \(x\) the same argument gives
\(\E_\nu[|x_T|\mid\mathrm{past}]\le |x|+\ell_j+2D_j\).

Write \(b,\ell,A,D\) for \(b_j,\ell_j,A_j,D_j\) in the current
stage. For \(w_i(x)=(x^2+\ell^2)^{-i/2}\), \(i=1,2\), let \(g_i\)
be the even function with \(g_i''=w_i\), \(g_i(0)=g_i'(0)=0\).
The logarithmic derivative of \(w_i\) is bounded by \(2/\ell\).
On a jump of size at most \(D=o(\ell)\), the quadratic Taylor
remainder for \(g_i\) is therefore
\((1+o(1))w_i(x)(\Delta x)^2/2\), uniformly.
Conditional means and the variance formula give
\begin{equation}\label{eq:lower-potential}
 \E_\nu\sum_{k<T} w_i(x_k)du_k
  =(2+o(1))\E_\nu g_i(x_T).
\end{equation}
The corresponding estimate starting from an intermediate
state subtracts \(g_i(x)\) on the right. One may first
sum to bounded stopping times and then pass to the limit
by boundedness of \(g_i(x_{k\wedge T})\) and nonnegativity of
the occupation sum. The two functions are
\[
 g_1(x)=|x|\operatorname{arsinh}(|x|/\ell)
                -\sqrt{x^2+\ell^2}+\ell,\qquad
 g_2(x)=(|x|/\ell)\arctan(|x|/\ell)
                -\tfrac12\log(1+x^2/\ell^2).
\]
Since \(g_1(b)=(1+o(1))bL\), \(g_2(b)\le Cb/\ell\), and
\(g_i(D)\le CD^2/\ell^i\), \eqref{eq:lower-potential} and the endpoint estimates
give \(\E_\nu\sum w_1du=(2+o(1))qbL\) and
\(\E_\nu\sum w_2du=O(1)\).
For the total resource \(R_j=\sqrt A\sum w_1du\) this proves
\begin{equation}\label{eq:lower-resource}
 \E_\nu[R_j\mid\mathrm{entry}]=(1+o(1))\mu_j,\qquad
 \mu_j=\sqrt2 B q\sqrt L.
\end{equation}
Also, conditional on any intermediate state \(x\), the
chance of eventual failure is at most \((|x|+\ell+2D)/b\).
Multiplying by \(ds=Aw_2(x)du\), summing predictable steps and
using \(|x|w_2(x)\le w_1(x)\) gives
\begin{equation}\label{eq:lower-failure-clock}
 \E_\nu[s_j\,1_{F_j}\mid\mathrm{entry}]
  \le \frac A b\,\E_\nu\sum\bigl(w_1(x)+(\ell+2D)w_2(x)\bigr)du
  \le (1+o(1))qS_j .
\end{equation}

We also need bounds from arbitrary intermediate states.
Convexity gives \(g_i(z)/z\) nondecreasing for \(z>0\),
so
\[
 \E_\nu[g_i(x_T)\mid\mathrm{past}]-g_i(x)
 \le \frac{g_i(b+D)}{b+D}(|x|+\ell+2D)-g_i(x).
\]
For \(i=1\) this is \(O(b)\), since
\(g_1(z)=|z|\log(1+|z|/\ell)+O(|z|+\ell)\),
\(|x|\log((2b)/(|x|+\ell))\le Cb\), and \(\ell L=o(b)\).
For \(i=2\), use
\(g_2(z)=\pi|z|/(2\ell)-\log(1+|z|/\ell)+O(1)\)
to bound it by \(O(L)\).
The conditional occupation bound thus gives
\(\E_\nu[R_{j,\mathrm{future}}\mid\mathrm{past}]\le Cb\sqrt A\)
and \(\E_\nu[s_{j,\mathrm{future}}\mid\mathrm{past}]\le CS_j\).
Writing \(R_j^2\) as the sum of the squared resource steps
plus twice the cross-products, and conditioning after the
earlier step, gives
\begin{equation}\label{eq:lower-resource-variance}
 \E_\nu[R_j^2\mid\mathrm{entry}]
 \le (r_j+Cb\sqrt A)\E_\nu[R_j\mid\mathrm{entry}]=O(q).
\end{equation}
For the clock tail, from
each active state the probability of accumulating another
\(K S_j\) units is at most \(1/2\) for a fixed sufficiently
large \(K\), by conditional Markov. Restarting at these
thresholds and using the per-step overshoot \(o(S_j)\)
gives the exponential tail. We cap each actual stage at
\(C S_j\log J\) with \(C\) large, so that the probability of
hitting any cap is \(O(J^{-5})\).

\subsection{Choice of the time allocation}

Let \(p_j^0=(1-q)^{j-1}\), and choose
\[
 T_0=\sum_{j=1}^J(p_j^0)^2\mu_j,\qquad
 \lambda_j=\frac{(1-\varepsilon)(p_j^0)^2}{T_0}.
\]
Then
\begin{equation}\label{eq:lower-resource-total}
 T_0/\sqrt L\ \longrightarrow\
 \frac{B}{\sqrt2}(1-e^{-2\gamma}),\qquad
 \lambda_j\asymp_{\gamma,B,\varepsilon}L^{-1/2}.
\end{equation}
Each resource step in stage \(j\) contributes time
\(\lambda_j r_j=1/n\). For the time comparison, continue the
ideal process after the first failure, with synthetic
restarts at zero, and append synthetic innovations if a
real branch is lost or budget-capped. Use the same tilted
single-step law in the appended part. It is defined for
all admissible states by \eqref{eq:lower-innovation} and \eqref{eq:lower-stage-clock}. The uniform
conditional estimates above hold for every appended stage.
The predictable means of \(\sum_j\lambda_j R_j\) sum to
\(1-\varepsilon+o(1)\) by \eqref{eq:lower-resource}, and the martingale part has
variance \(O(\sum_j\lambda_j^2 q)=O(1/L)\) by \eqref{eq:lower-resource-variance}.
Thus the chance of exhausting the budget is \(O(1/L)\).
Similarly the ideal reward \(\sum_j\sqrt{\lambda_j}R_j\)
has second moment \(O(\sqrt L)\) because its sum of conditional
means is \(O(L^{1/4})\), and its martingale part has
variance \(O(\sum_j\lambda_j q)=O(L^{-1/2})\).

Before the first failure, the probability of entering stage
\(j\) is \((1+o(1))p_j^0\) uniformly, since the relative error
in every conditional failure probability is \(o(1)\) and
\(Jq=\gamma\). The expected ideal reward up to and including
the first failure is therefore
\((1+o(1))\sqrt{(1-\varepsilon)T_0}\).
If the retiling loss probability summed over the branch
is \(O(m^{-9})\), as arranged below, all three types of
artificial truncation have probability \(o(1)\).
Cauchy--Schwarz and the reward second moment show that they
change the mean reward by \(o(L^{1/4})\). Converting from
reward to the sum of the spatial block coefficients gives
\begin{equation}\label{eq:state-reward-sharp}
 \E_\nu\sum_{\rm retained\ steps}\sqrt{r_j}
 \ge (1-o(1))\sqrt n\,\sqrt{(1-\varepsilon)T_0},
\end{equation}
because \(r_j^{-1/2}=\sqrt{n\lambda_j}\).

For the terminal cost, the conditional \(\nu\)-mean of
\(\log(1+\beta)\) equals \(ds/2+O(ds\,\sup|\beta|)\)
by expanding \((1+\beta)\log(1+\beta)\) under the ordinary
law. The conditional variance is \(O(ds)\), and the
centered jumps are \(O(m^{-c_0})\). In stage \(j\) the
total actual clock, including previous stages, is bounded
by \(C S_j\log J\). The centered sum of the log increments
up through each such stage, stopped at its caps, is at most
\(m^{-\kappa/4}S_j\) in absolute value except with probability
\(\exp(-m^{c_1})\) for some \(c_1>0\), by the bounded-increment
martingale tail inequality. To obtain it, apply the same
conditional exponential bound as for bounded independent sums.
If \(d_i\) are the centered martingale increments, then for
\(|t|\max_i\|d_i\|_\infty\le1\),
\(\exp(t\sum_{i\le j}d_i-Ct^2\sum_{i\le j}
\E[d_i^2\mid\mathcal F_{i-1}])\) is a supermartingale.
Stop it at the first threshold crossing and optimize \(t\),
using both signs. Here the predictable quadratic variation
bound is \(C S_j\log J\), and
\(m^{-\kappa/2}S_j/\log J\ge m^{\kappa/2}/\log J\).
Union over \(j\) preserves an exponential bound, and the
accumulated drift errors are \(o(S_j)\). Also the clock of
earlier stages is \(O(S_j L^{-4}\log J)=o(S_j)\).
Thus on one joint good event, in stage \(j\) one has
\(\log Y=s_j/2+o(S_j)\), with \(s_j\) the clock spent in
that stage.

A completed success has relative ratio at most \(D_j\).
On failure in stage \(j\) the actual ratio satisfies
\(|M|\le (1+o(1))b_j\), including all prior overshoots.
Apply Cauchy--Schwarz with \eqref{eq:lower-failure-clock} and \(p_j=(1+o(1))q\).
Conditional on ideal entry, the expectation on actual
failure and on the joint good clock event is bounded by
\[
 (1+o(1))b_j\sqrt{(1+o(1))q}\,
       \sqrt{(1+o(1))q S_j/2+o(qS_j)}
 = (1+o(1))qB/\sqrt2.
\]
This uses a bound on the integrand on the joint good
event only, so no conditional estimate on the previous
value of \(\log Y\) at a rare entry history is needed.
The terminal cost on time, geometry or clock stops is
\(O(B\sqrt{\log J})\) on this same event, with expected
contribution \(o(1)\) from the above bounds. On completion
of all stages, \(|M|\le\sum_j D_j=o(b_J/\sqrt{\log J})\)
gives another \(o(1)\) cost. On the bad clock event the
deterministic bound
\(|\log Y|\le C\sqrt{n\log R\, S_J\log J}\)
follows by Cauchy--Schwarz on the increments. Its
contribution is negligible because the event has
exponentially small probability. Summing the entry
probabilities times the failure costs gives
\begin{equation}\label{eq:state-cost-sharp}
 \E_\nu[|M|\sqrt{(\log Y)_+}]
 \le (1+o(1))\frac B{\sqrt2}(1-e^{-\gamma}).
\end{equation}
The same bounds give
\(\E_\nu[M^2\log(2+Y|M|/\eta_*)]=O(L)\) for any fixed
\(\eta_*>0\). Indeed on the joint good event the
logarithm is bounded by \(C(1+S_j\log J)\) in the
stopping stage and \(|M|\le C b_j\). On its complement
the preceding deterministic bound again suffices.

\subsection{Realizing the tilted process by a test potential}

Begin with \(u_0(x)=x_2-1/2\) on \(\Omega\).
At a step, a branch cell is an old affine triangle (or the
initial square), and its gradient and stage parameters are
constant. Choose \(A\) by
\[
 A^{-\mathsf T}
   =\begin{pmatrix}a_j^{-1/4}&M a_j^{1/4}\\0&a_j^{1/4}\end{pmatrix}.
\]
This has determinant one. Both \(A\) and \(A^{-1}\) have
norm bounded by fixed powers of \(m\) since \(a_j\) and \(a_j^{-1}\)
do by \eqref{eq:lower-stage-clock} and \(|x|\le b_j\), and \(|M|\le 2b_1\).
The same polynomial bound holds for the change between
successive charts. In the new chart retain all grid squares
inside the old cell with side the old intrinsic leg times a
dyadic factor comparable to \(m^{-D_0}\), for a sufficiently
large fixed \(D_0\). The old intrinsic leg of a triangle is
the side of the square in which it was one of the four
reference triangles. An uncovered point lies within a
constant times the new side of the old boundary. Because
the aspect ratio in the new chart is polynomial, the ratio
of boundary-strip area to old area is \(O(m^{-10})\) for
large enough \(D_0\). Freeze the potential on the uncovered
remainder. By the pathwise tilt identity the lost tilted
probability over \(O(m)\) steps is \(O(m^{-9})\).

On each new chart square \(Q=x_Q+dAS\), write \(p_Q\) for
the pullback of a reference function \(p\) by
\(x\mapsto d^{-1}A^{-1}(x-x_Q)\), extended by zero.
Reveal the finitely many Gaussian coefficients
\(d^{-1}W(g_{Q,k})\), where \(g_{Q,k}=(g_k)_Q\).
Build \(U\) on \(S\) from these reference coefficients and
add \(dY\sqrt{r_j}U_Q\) to the potential. Its physical
gradient increment is the specified matrix applied to
\(\xi\), by the chain rule. Its mean and integrated
gradient are zero. Its support is strictly inside \(Q\).
The potential therefore remains globally continuous
piecewise affine, hence Lipschitz for each finite mesh,
and its vertical derivative stays positive since
\(\sup|\beta|=o(1)\).

Each reveal is predictable. For a new square, every
previous direction is either zero there or affine there.
This follows from the nested placement inside previous
child triangles, since all the old \(g\)'s are piecewise
affine on those triangles. The new directions have zero
affine moments. They are orthogonal to all previous
revealed directions and to the constant function.
The directions in disjoint new squares are orthogonal
as well. To justify conditional freshness, at each
step orthonormalize the new finite family given the
old data. Conditional on that data, the remaining white
noise is isonormal on the orthogonal complement of the
revealed space. Conditioning further on the new independent
coordinates leaves it isonormal on the new complement.
This induction starts with the ordinary white noise on
the square and establishes the claimed reference law
and all the conditional estimates for the chosen charts.

A microstep costs at most
\(\log R+D_0\log m+O(1)\) of intrinsic log-depth.
Choose \(n\) to be the largest integer whose cumulative
bound is at most \((1-\epsilon_m)m\), with
\(\epsilon_m=(\log m)^{-1/2}\).
Then \(n\sim m/\log R\) and every used chart leg is at
least \(d_{\min}=h^{1-\epsilon_m}\) up to an absolute
factor. This is enough to pair against \(P_hW^0\).
Conditional on the old data, the normalized scalar
pairing of a block differs from its unprojected expected
pairing only by \(o(1)\) before division by \(\sqrt v\).
Indeed, the sum of the unclipped diagonal covariance
differences is bounded by
\[
 C\sum_k \lambda_k^{-1}(h\|A^{-1}\|/d)\|\nabla f_k\|_2
 \le C(h\|A^{-1}\|/d)K_{\max}=o(1).
\]
The estimate uses cell Poincar\'e for \(P_h\) and the
zero extension of \(f_k\), which is supported inside \(Q\).
For a predictable adaptive test function \(P_h f_{Q,k}\),
write its noise pairing as its conditional old component
plus the remaining centered Gaussian pairing. The old
component makes zero contribution in expectation since
the fresh input coefficients are odd and the clipping event
is even. The remaining component has variance no greater
than \(\|P_h f_{Q,k}\|_2^2\le C d^2\).
Fourth moments and the clipping probability then bound the
normalized pairing loss by \(O(R^{-5}\log R)\) before
division by \(\sqrt v\). Thus the conditional expected
contribution to \(W^0(P_hu)\) from \(Q\) is
\[
 (1+o(1))d^2Y\sqrt{r_j}\sqrt{\log R/\pi},
\]
uniformly in the old history. The factors \(d^2Y\) are
exactly the probabilities of the squares under the tilted
spatial branch law. The initial affine potential has
expected pairing zero. Sum over the retained steps and use
\eqref{eq:state-reward-sharp} to obtain
\[
 \E W^0(P_hu)\ \ge\
 (1-o(1))\sqrt{m/\pi}\,\sqrt{(1-\varepsilon)T_0}.
\]

\subsection{Duality and optimization}

Let \(D_*(u)=1+\eta_*+\int_\Omega
\eta_*\Theta^*(|\partial_1 u|/\eta_*)\) for fixed
\(\eta_*>0\), where \(\Theta^*\) is the convex conjugate of
\(\Theta(t)=e^{t^2}-1\). If a current has gauge at most \(r\),
apply the weak divergence identity against \(u\), use
Fenchel in the first coordinate, and bound the second
coordinate by \(r\int_\Omega|\partial_2u|=r\). This gives
\(|W^0(P_hu)|\le r D_*(u)\). To use a Lipschitz test in the
weak identity, extend it by reflection, smooth, and pass
to the limit in the integrals against the integrable
current using bounded gradients and almost everywhere
convergence. Infimizing over the currents gives the same
bound with \(r=\mathcal F(P_hW^0)\).

For \(v\ge0\), the maximizing \(t\ge0\) in \(\Theta^*(v)\)
satisfies \(v=2te^{t^2}\), so
\(\Theta^*(v)\le v\sqrt{\log(2+v)}\). Indeed the bound on
\(t\) is immediate for \(t\le\sqrt{\log 2}\), and otherwise
\(2t>1\) implies \(t^2\le\log v\). Since \(X=YM\) and
\(|M|\le2b_1=o(1)\), the integral in \(D_*\) is at most
the conditional spatial \(Y\)-average of
\(|M|\sqrt{(\log Y)_+}+o(1)\). For example,
\(\log(2+Y|M|/\eta_*)\le(\log Y)_++\log(2+2b_1/\eta_*)\).
Equation \eqref{eq:state-cost-sharp} and Jensen under
the probability density \(Y\) now give
\[
 \E D_*(u)\le 1+\eta_*+(1+o(1))\frac B{\sqrt2}(1-e^{-\gamma}),
 \qquad \E D_*(u)^2=O(L).
\]
The second bound uses the terminal estimate for
\(\E_\nu[M^2\log(2+Y|M|/\eta_*)]\) above.

The grid cut estimate \eqref{eq:sc-box-cut} and affine
interpolation bound the flux gauge of the histogram with cell
values \(g_v/h\) by \(C\|g\|_2\) for mean-zero \(g\).
Here we use \(\|J_1\|_{\psi_2}\le \|J_1\|_\infty/\sqrt{\log 2}\).
The gauge is subadditive by the triangle inequality, so it is
uniformly Lipschitz in the Gaussian cell coefficients.
Gaussian concentration gives \(\Var(\mathcal F(P_hW^0))=O(1)\). Cauchy--Schwarz
therefore bounds the covariance of the gauge with \(D_*\)
by \(O(\sqrt L)\), negligible in the pairing bound.
With \eqref{eq:lower-resource-total} and \(L\sim\log\log m\), we obtain
\[
 \liminf_{h\downarrow0}
 \frac{\E\mathcal F(P_hW^0)}{\sqrt m(\log\log m)^{1/4}}
 \ \ge\
 \frac1{\sqrt\pi}\,
 \frac{\sqrt{(1-\varepsilon)(B/\sqrt2)(1-e^{-2\gamma})}}
      {1+\eta_*+(B/\sqrt2)(1-e^{-\gamma})}.
\]
Set \(B=\sqrt2/(1-e^{-\gamma})\), then let
\(\varepsilon,\eta_*,\gamma\) tend to zero after the
mesh limit. The right side tends to \(1/\sqrt{2\pi}\).
The variance bound also gives the lower bound with this
constant in probability.

\section{The sharp upper bound}
\label{sec:sharp-upper}

The upper construction uses a thick Gaussian refinement step in
a determinant-one chart. Its principal singularity gives the
variance coefficient \(1/(4\pi)\). The following phase and
stability estimates retain that coefficient through the
predictable chart changes.
\subsection{A sharp continuous phase}\label{sec:sharp-phase}
Let \(M=\log(1/h)\) and, throughout this subsection, write
\(a=\log\log M+10\log\log\log M\sim\log\log M\), \(a_0=a^{1/10}\),
\(u_0=a^{-1/10}\), \(L_0=a^{1/4}\), and \(K=a^4\). Choose a smooth
\(\alpha(u)\), constant \(a_0\) for \(u\le a^2\), increasing to
\(a\) over a window of length \(O(a)\) in \(\log u\), and constant
afterwards. Such a ramp is obtained by integrating a smooth bump on
an interval of length \(C a\), with \(C\) a large fixed constant.
Its first derivative in \(\log u\) is in \([0,1/8]\), and the
higher derivatives used below are bounded. Put
\[
 r=e^{-\alpha(u)}\sqrt{\frac{u}{1+u}}(1+u/M)^{-1/4},
 \qquad I(u)=K\int_0^u r(v)\frac{dv}{v}.
\]
For \(u\le1\), \(r(u)\le e^{-a_0}\sqrt u\). From \(1\) to
the end of the ramp the logarithmic interval has length \(O(a)\)
and \(r\le e^{-a_0}\). Beyond that point, integration up to \(M\)
costs at most \(e^{-a}\log M\), and the factor \((1+u/M)^{-1/4}\)
bounds the remaining integral by \(C e^{-a}\). Thus
\(I(\infty)=o(1)\). Let
\[
 h(t,y)=\frac{L(t)^2-y^2}{L_0^2},\qquad
 L(t)=L_0+2(1+\delta_M)\sqrt a\,t,
\]
where \(\delta_M\to0\) will dominate errors in the derivative bounds.
Use the implicit phase
\[
 -h(t,y)=I(u)-r(u)(e^{\alpha(u)x^2/u}-1)
\]
on the strip \(h+I(\infty)>0\), setting \(u=0\) on the central
line where no positive solution exists. The right side is strictly
increasing in \(u\). Write \(s=x/\sqrt u\), \(z=\alpha s^2\),
\(J=e^z\), \(\beta=\partial_{\log u}\log r\),
\(\gamma=\partial_{\log u}\log\alpha\), \(d=1-\gamma\), and
\[
 p=2\alpha s rJ,\quad q=2\alpha rJ(1+2z),\quad
 Q=r[K-\beta(J-1)+dzJ],\quad
 D=q(Q+sQ_s/2-Q_{\log u})-Q_s^2
\]
where \(Q_{\log u}\) is at fixed \(s\).
We have \(\beta\in[-3/8,1/2]\), \(|\beta'|\lesssim1\), and
\(|\gamma|+|\gamma'|\lesssim 1/a_0\). Hence \(Q\asymp r(K+zJ)\).
The bounds below hold uniformly in \(u,z\).
\[
 \frac qQ\le4\alpha(1+o(1)),\qquad
 \frac{p^2}Q\le4\alpha(1+o(1))rJ,\qquad
 0\le\frac{D}{Q^2}\lesssim\alpha,\qquad
 \frac{u\widehat H_{xx}}{\widehat H_Y}\lesssim
     \alpha rJ(1+\log K).
\]
Here \(\widehat H(x,Y)\) is the implicit level for the graph in
\((x,Y)\), so \(H(t,x,y)=\widehat H(x,-h(t,y))\).
For the first two bounds, split at \(z=(\log K)/2\). Below this
value the numerator relative to the \(rK\) term is negligible.
Above it, use \(dz-\beta\ge(1-o(1))z\).
To check the other two, put \(k=K+\beta\), \(v=dz-\beta\),
\(w=v+d\), \(A=(1-\beta)k-\beta'\), and
\(E=v(1+v)+d^2z-d'z+\beta'\). Then
\[
 Q=r(k+Jv),\quad Q_s=wp,\quad Q+sQ_s/2-Q_{\log u}=r(A+JE)
\]
and
\[
 \frac{D}{2\alpha r^2J}=(1+2z)A+J\big[(1+2z)E-2zw^2\big].
\]
The bracket is a quadratic in \(z\) with uniformly bounded
coefficients, leading coefficient \(2d-d^2-2d'\ge1/2\).
For bounded \(z\) its negative part is absorbed by \(A\asymp K\);
for large \(z\) it is positive. Together with \(q>0\), this makes the implicit Hessian positive
semidefinite. Indeed, in the coordinates \((x,u)\) the quadratic
form for \(D^2\widehat H/\widehat H_Y\) has diagonal entries
\(q/u\) and \((Q+sQ_s/2-Q_{\log u})/u^2\), off-diagonal
entry \(-Q_s/u^{3/2}\), and determinant \(D/u^3\). Also
\[
 \frac{u\widehat H_{xx}}{2\alpha rJ\widehat H_Y}
 =(1+2z)-\frac{4zJw}{k+Jv}
           +\frac{2zJ(A+JE)}{(k+Jv)^2}.
\]
For \(z\le1+\log K\) use \(Q\asymp r(K+zJ)\). For larger \(z\),
replace \(k+Jv\) by \(Jv\) at error \(O(K/J)\), and use
\(1-2w/v+E/v^2=O(1/z)\). This gives the stated bound.

Let \(B_H=D^2H+\nabla H\nabla H^{\mathsf T}\).
For positive levels the graph identities give
\[
 \frac{\det B_H}{\widehat H_Y^2}
 = h_y^2\left(\frac qQ+\frac{D}{uQ^2}\right)
  +(-h_{yy})\left(\frac{\widehat H_{xx}}{\widehat H_Y}+\frac{p^2}Q\right),
 \qquad -H_t=h_t\widehat H_Y.
\]
For \(u\ge a^2\), the displayed derivative bounds therefore give
\[
 \frac{\sqrt{\det B_H}}{\widehat H_Y}
 \le 2\sqrt a(1+o(1))
        \sqrt{h_y^2-h_{yy}(h+I+r)}
 \le 4\sqrt a(1+o(1))L/L_0^2.
\]
Indeed \(rJ=h+I+r\), and the last bound uses
\(|y|^2\le L^2+L_0^2 I(\infty)\).
For \(u_0\le u\le a^2\) the same inequality holds with a
strict \(o(\sqrt a)\) coefficient, since \(\alpha=a_0\)
and \(a_0(1+\log K)/u_0=o(a)\).

Flatten \(H\) by replacing it with \(\chi(H)\), where
\(\chi'=0\) on \([0,u_0/2]\), \(\chi'=1\) on \([u_0,\infty)\),
and \(\chi'\) increases smoothly between these regions with
\[
 \chi''\le C u_0^{-1}\chi'(1+\log^2(1/\chi')).
\]
Such a cutoff follows by starting its derivative like \(e^{-1/v}\)
in the rescaled transition coordinate. In the transition with
\(\chi'\ge e^{-4a_0}\), the determinant of
\(B_{\chi(H)}=\chi' D^2 H+(\chi''+\chi'^2)\nabla H\nabla H^{\mathsf T}\),
divided by \((\chi'\widehat H_Y)^2\), gains a factor
\(\chi''/\chi'+\chi'\) only on the gradient rank-one term.
This follows by the same graph formula with \(q/Q\) and \(p^2/Q\)
multiplied by that factor. Its square root is still bounded by
\(o(\sqrt a)\sqrt{h_y^2-h_{yy}(h+I+r)}\), since
\(a_0^3/u_0=o(a)\).
In the remaining transition, the two coefficients \(\chi'\) and
\(\chi''\) are at most \(e^{-4a_0}\operatorname{poly}(a)\)
and \(\widehat H_Y=u/Q\le e^{a_0}\operatorname{poly}(a)\).
The determinant root is then \(o(1)\) absolutely.
Set \(F=\chi(H)+a^{-1}(T-t)\), with \(T=1+o(1)\).
Choose \(\delta_M\downarrow0\) to dominate the relative
errors in the high-level bounds, for example a sufficiently
large multiple of \((\log a)^{-1/4}\).
The low-level bound \(e^{-3a_0}\operatorname{poly}(a)=o(a^{-1})\)
is absorbed by the time-only term.
For \(B_0=D^2F+\nabla F\nabla F^{\mathsf T}\), use the
regularized matrix
\[
 B=B_0+\frac{\epsilon_{\rm reg}^2}{1+\operatorname{tr}B_0}I,
 \qquad \epsilon_{\rm reg}=(\log M)^{-1}.
\]
The increase in determinant root is \(O(\epsilon_{\rm reg})\).
The bounds just proved give
\(-F_t\ge(1+\omega_M)\sqrt{\det B}\) for a slowly vanishing
\(\omega_M>0\) after choosing the wall-speed margin as above.
They are uniform over \(0\le t\le T\).



On \(0\le H\le M\) the same level equation gives
\[
 x^2\le u+\frac{u}{\alpha(u)}
       \left[O(\log a)+\tfrac12\log(1+1/u)
                                  +\tfrac14\log(1+u/M)\right]
       \le (1+o(1))u+o(1).
\]
Here \(\alpha\ge a_0\), \(u\log(1+1/u)\le1\), and
\(\chi(u)=u+O(u_0)\). Hence \(x^2\le(1+o(1))\chi(H)+o(1)\)
on the active states and their retained one-step overshoots.
The strip has vertical radius
\(L_0+2(1+o(1))\sqrt a\,T+o(1)\). These estimates, together
with the determinant bound, are the phase inputs used below.

\subsection{Discrete stability of the sharp continuous phase}
\label{sec:sharp-discrete-stability}

We prove the pointwise conditional estimate that allows the sharp
continuous phase to run along the adaptive spatial construction.
All vectors in this subsection are in the unbalanced coordinates
of the upper bound. The preliminary determinant-one balancing map
takes the initial domain to a rectangle tiled by squares in these
coordinates. The initial means on these squares are handled by a
separate current. The field entering the phase starts at zero.

Use \(M=m\) in Section~\ref{sec:sharp-phase}. Denote the deterministic
time horizon by \(T_{\mathrm{end}}=1+o(1)\) and the normalized
variance allowance per split by \(\delta\asymp(\log m)^2/m\).
Write \(u=H(t,z)\), \(v=\nabla_z F(t,z)\), and
\[
 B_0=D_z^2F+vv^{\mathsf T},\qquad
 \eta=\frac{\epsilon_{\mathrm{reg}}^2}{1+\tr B_0},\qquad
 B=B_0+\eta I,\qquad K=\sqrt{\det B},\qquad C=KB^{-1}.
\]
The selector \(C\) is positive definite with determinant one.
There is a slowly vanishing \(\omega=\omega_m>0\) for which
\(-F_t\ge(1+\omega)K\), and \(F(0,0)=o(1)\). We use the deterministic
level cutoff \(U_*=m/(\log m)^A\), with \(A\) a large fixed constant.

\paragraph{Bounds for the selector.}
In the region \(u\le 2U_*\) use widths \(b_x=\sqrt{1+u}\) and
\(b_y=1\). In statements about derivatives one can smoothly
replace \(u\) by zero for \(u\le u_0/4\), and leave it as \(u\) for
\(u\ge u_0/2\). This does not affect any of the estimates for \(F\),
which is spatially flat for \(u\le u_0/2\).
For some fixed power \(P\) of \(\log m\), independent of \(A\),
\begin{equation}\label{eq:sharp-width-derivatives}
 \frac{1+u}{P}\le K\le P(1+u),\qquad
 C_{ii}\le P b_i^2,\qquad
 |b_i b_j F_{ij}|\le P(1+u).
\end{equation}
The last bound holds also after applying \(\partial_t\) or one
of \(b_k\partial_k\) to \(F_{ij}\) before multiplying by
\(b_i b_j\). These operators applied to
\(\log(1+u)\) and \(\log K\) are bounded by \(P\).
Here and below \(P\) can be increased by a fixed power.

To verify these bounds, let \(Y=-h(t,y)\), \(H_Y=u/Q\),
and \(v_H=\nabla H\). Use the profile quantities \(p,q,Q,D,s\)
of Section~\ref{sec:sharp-phase}, with the profile constant
there denoted by \(K_r=a^4\). Differentiating the inverse graph
gives the exact positive semidefinite decomposition
\begin{equation}\label{eq:sharp-hessian-factors}
 D^2H=H_Y\left[
 \frac q u\left(e_x-\frac{Q_s}{q\sqrt u}v_H\right)
             \left(e_x-\frac{Q_s}{q\sqrt u}v_H\right)^{\mathsf T}
 +\frac{D}{q u^2}v_Hv_H^{\mathsf T}+Y_{yy}e_ye_y^{\mathsf T}\right].
\end{equation}
Indeed, for the function \(Y=G(u,x)\) defining the graph,
\(G_{xx}=-q/u\), \(G_{xu}=Q_s/u^{3/2}\), and
\(G_{uu}=(Q_{\log u}-sQ_s/2-Q)/u^2\). Implicit differentiation
and completing the square give \eqref{eq:sharp-hessian-factors}.
Also \(H_x=\sqrt{u}\,p/Q\), \(H_y=uY_y/Q\) and \(H_t=uY_t/Q\).

For \(u\ge u_0/4\), \(r^{-1}\le(\log m)^{O(1)}\) because
\(\alpha\le a\) and \(e^a=(\log m)(\log\log m)^{10}\).
On the strip \(rJ=h+I+r=O(\sqrt a)\). Thus \(J\), \(s\), \(z\),
\(Q^{-1}\), \(q^{-1}\), \(Y_y\), \(Y_{yy}\), and all the numerators
in the graph formulas are bounded by powers of \(\log m\).
The formula for \(D\) gives
\[
 D\gtrsim \alpha r^2J\big[(1+z)K_r+(1+z^2)J\big].
\]
For bounded \(z\) the term with \(K_r\) absorbs the bracket's
bounded negative part. For sufficiently large \(z\) the bracket
has its uniform positive quadratic lower bound.
Using the three first derivatives of \(H\) just computed shows
that \(\log u\) and \(s=x/\sqrt u\) have derivatives bounded by
powers of \(\log m\) in the scaled directions and in time.
The ramp of \(\alpha\) has bounded derivatives through order four
in \(\log u\), so the same bounds apply to the derivatives of
\(q,Q,D\) and to their logarithmic derivatives. These assertions
follow by differentiating their explicit formulas. The displayed
positive lower bounds for \(q,Q,D\) ensure no larger reciprocal
appears.

The three vector factors of \eqref{eq:sharp-hessian-factors},
with the square roots of their scalar coefficients included,
have coordinates bounded by \(P(1,\sqrt{1+u})\), together
with their scaled derivatives and their time derivatives.
For example the first scalar coefficient is \(q/Q\), and
\(Q_s/(q\sqrt{u})\) and its derivatives have size at most
\(P/\sqrt{1+u}\). The second scalar coefficient is \(D/(q uQ)\)
and has square root at most \(P/\sqrt{1+u}\). The third
factor is \(\sqrt{uY_{yy}/Q}\,e_y\).
Meanwhile \(v_H\) and its corresponding derivatives are bounded
coordinatewise by \(P(\sqrt{1+u},1+u)\).
The cutoff \(\chi\) can be chosen so the square roots of
\(\chi'\) and \(\chi''\) have smooth, polynomially bounded
derivatives in \(1/u_0\). For example, choose \(\chi''\) as a
normalized positive flat exponential bump on the transition
interval, then integrate twice with zero initial data.
This proves the Hessian estimates for \(F=\chi(H)+
\epsilon_0(T_{\mathrm{end}}-t)\).

For \(u\ge u_0\) the determinant identity in the continuous
phase gives \(\det B_0=(u/Q)^2\mathcal D\), where
\(\mathcal D\) and its derivatives in the stated directions are
polynomially bounded in \(\log m\). Also \(\mathcal D\ge P^{-1}\).
When \(|y|\ge L(t)/2\), use \(h_y^2q/Q\). Otherwise \(h\) is
bounded below by an absolute positive constant, so \(rJ=h+I+r\)
is bounded below and \(z=\log J\) is bounded away from zero.
Then use \(Y_{yy}p^2/Q\). The change from adding \(\eta I\)
in the determinant is
\(\epsilon_{\rm reg}^2\tr B_0/(1+\tr B_0)+\eta^2\).
Its scaled and time derivatives are polynomially bounded, since
\(\tr B_0\gtrsim (1+u)/P\) and its derivatives are at most
\(P(1+u)^2\). In the cutoff layer all entries and derivatives
are polynomially bounded and
\(\det B\ge\epsilon_{\rm reg}^4/(1+\tr B_0)^2\).
This gives the bounds for \(K,\log K\).
The cofactor formula for \(C\) and the coordinate bounds on the
Hessian factors and gradient give \(C_{ii}\le P b_i^2\).

For a quantitative estimate when the selector moves, let \(S_C\)
have rows given by the three Hessian factors in
\eqref{eq:sharp-hessian-factors} multiplied by \(\sqrt{\chi'}\),
the two vectors \(\sqrt{\chi''}v_H\) and \(\chi'v_H\), and the
two identity rows multiplied by \(\sqrt{\eta}\), all divided
by \(\sqrt K\). Thus \(S_C^{\mathsf T}S_C=C^{-1}\).
For any fixed determinant-one metric \(G\) with
\(G_{ii}\le P_1 b_i^2\), where \(P_1\) is a fixed power of \(\log m\),
the derivatives of \(S_CG^{1/2}\) in the same scaled and time
directions are at most \(P_2\sqrt{1+u}\). For each row this follows
by the coordinate bounds and \(\|e_i^{\mathsf T}G^{1/2}\|
=\sqrt{G_{ii}}\). For the regularization rows use
\[
 |\partial\sqrt{\eta}|\le
 \frac{\epsilon_{\rm reg}|\partial\tr B_0|}
 {2(1+\tr B_0)^{3/2}}
 \le P\sqrt{1+u}
\]
outside the cutoff layer and polynomial bounds inside it.

\paragraph{The induction on active cells.}
At generation \(j\), \(t=j\delta\) and the accumulated regular
field is \(Z_j\). The metric \(G\) of a parent is the product of
its chart map with its transpose. Its determinant is one.
All root cells have \(G=I\), \(Z_0=0\) and selector \(C=I\).
Select an interior representative \(x_Q\) of each active cell
measurably from the old history. All pointwise statements
are on cell interiors, with cell boundaries omitted from
spatial integrals. Freeze terminal cells and cells whose representative
has \(H(t,Z_j(x_Q))\ge U_*/2\).
At each other parent set \(C_Q=C(t,Z_j(x_Q))\), diagonalize
\(G^{-1/2}C_QG^{-1/2}\), clip its two reciprocal eigenvalues
to \([K_0^{-1},K_0]\), and conjugate back by \(G^{1/2}\).
The result \(G_+\) satisfies
\begin{equation}\label{eq:sharp-metric-clip}
 \det G_+=1,\qquad K_0^{-1}G\preceq G_+\preceq K_0G,\qquad
 G_+\preceq K_0^{-1}G+C_Q .
\end{equation}
The last inequality follows from
\(\operatorname{clip}(\lambda)\le K_0^{-1}+\lambda\)
for each eigenvalue.

We use the following properties of the two local splits in the
assembly. At fixed \(x\), the uncut increment is Gaussian with
covariance at most \(\delta G_+\). Its coarse and fine summands
are independent conditional on the old history, and each is
multiplied by an even scalar in \([0,1]\) depending only on
its own Gaussian data. There is a nonnegative retained envelope
\(E_j(x)\) with \(E_j(x)\le C E_0\), \(E_0=\sqrt m/(\log m)^B\),
and with conditional tail \(C e^{-c v^2}\) at \(v\).
The retained increment has chart magnitude at most
\(C\sqrt\delta E_j(x)\) and a Hölder bound with a fixed exponent
and the same coefficient on the child containing \(x\),
with distances divided by its chart leg. For the coarse
increment this follows from its bound on the containing base
cell. For the fine increment it follows from its bound on
the fine child. A rejected fine base cell contributes zero
throughout that cell. The assembly and the local estimates
verify these assertions before either reveal at this step.
The child contraction in chart diameter, combined with the
fixed relative distortion in \eqref{eq:sharp-metric-clip},
can be taken as small as desired for descendants that split
again.

Define the past envelope sum \(T_j(x)\) by \(T_0=0\) and
\(T_{j+1}=\rho T_j+E_j^2\), with fixed \(0<\rho<1\).
Along an active history assume the relative width changes
at each past step are at most \(1+o(1)\), and maintain
\(G_{ii}\le P_1 b_i(t,Z_j(x))^2\) at every point in the cell.
For \(x,x'\) in one old cell the ancestor Hölder bounds imply
\begin{equation}\label{eq:sharp-discrete-oscillation}
 \sum_i\frac{|Z_j(x)_i-Z_j(x')_i|}{b_i(t,Z_j(x))}
 \le P_2\sqrt\delta\,\sqrt{T_j(x)} .
\end{equation}
To see this, the distance between \(x\) and \(x'\) in the chart
of generation \(i<j\), divided by the leg of their common child
at that generation, is at most \(C\vartheta^{j-i-1}\) for a fixed
small \(\vartheta\). Apply the ancestor Hölder estimate there.
On converting that chart increment to coordinate \(k\),
the factor is at most \(\sqrt{P_1}b_k(t_i,Z_i(x))\).
Its ratio to the current width is bounded by
\(C(1+c)^{j-i}\) for a small fixed \(c\).
Choose \(\rho\) larger than the square of the resulting
geometric decay ratio and use Cauchy--Schwarz. This gives
\eqref{eq:sharp-discrete-oscillation}, and \(T_j\lesssim E_0^2\).

For clarity, the deterministic induction starts with
\(T_0=0\) and all widths equal to one. Immediately after any
split the old states and their new values, including the
time advance, have level at most \(2U_*\), as follows.
Use the scaled derivative bounds above along the short
segments in state and time. A state displacement at most
\(C P_2\sqrt\delta E_0\) in the widths changes \(1+u\) and
\(K\) by a relative \(o(\omega)\), and a time change \(\delta\)
has the same property. This follows by a continuation
argument on the first segment where the width ratio lies
in \([1/2,2]\), because the logarithmic derivatives are
bounded by \(P\). The segments remain in the strip.
At positive levels its margin in the variable \(Y\) is
\[
 I(\infty)-I(u)+r(J-1)\ \ge\ I(\infty)-I(u)
 \ \gtrsim K_r r(u),
\]
by integrating on \([u,2u]\). This is at least an inverse power
of \(\log m\) for \(u_0/4\le u\le 2U_*\).
At lower levels the margin is bounded below by
\(I(\infty)-I(u_0/4)\). The same estimates and the explicit
polynomial bounds for \(Y_t,Y_y\) handle the short change
of time and the vertical coordinate. Take \(B\) sufficiently
large for these estimates.

In particular, if the representative has level below \(U_*/2\),
\eqref{eq:sharp-discrete-oscillation}, applied with the
representative as the first point, shows that all old
states in that cell have level below \(U_*\). In this
application the derivative continuation begins at the
representative, so it requires no pointwise upper-level
assumption beyond the induction for the preceding step.
The pointwise widths at the representative and any point
of the old cell are relatively \(1+o(1)\).
By \eqref{eq:sharp-metric-clip} and \eqref{eq:sharp-width-derivatives},
\(G_{+,ii}\le (1-c_1)P_1 b_i^2\), for some fixed \(c_1>0\)
and \(P_1\) sufficiently larger than \(P\).
The fresh retained increment consequently has normalized
width displacement at most \(C\sqrt{P_1\delta}E_j(x)\).
The same continuation estimate keeps the new level below
\(2U_*\) and closes the metric cap and width induction.
Geometric terminal cells reached by a split already satisfy
these estimates and are frozen thereafter. If a representative
triggers the high-state stop, \eqref{eq:sharp-discrete-oscillation}
and the same continuation estimate show that the entire cell
has level at least \(U_*/3\), for large \(m\).

\paragraph{Lag of the implemented metric.}
For positive determinant-one matrices set
\(\mathcal R(G,C)=\tr(GC^{-1})-2\).
If the larger relative eigenvalue is \(e^d\), then
\(\sqrt{\mathcal R(G,C)}=2\sinh(d/2)\).
For clipping at a fixed target, \(d\) becomes
\((d-\log K_0)_+\), and
\(\sinh((d-\log K_0)/2)\le K_0^{-1/2}\sinh(d/2)\)
when \(d\ge\log K_0\). Thus clipping contracts the
square-root lag by \(K_0^{-1/2}\) without an additive term.
Also
\[
 \sqrt{\mathcal R(G,C)}=\sigma-\sigma^{-1},
 \qquad \sigma=\|S_CG^{1/2}\|_{\mathrm{op}}\ge1.
\]
This expression is \(2\)-Lipschitz in the rectangular matrix
factor, since operator norm is \(1\)-Lipschitz and
\(s-s^{-1}\) has derivative at most two on \([1,\infty)\).

Write \(\mathcal R=\mathcal R(G,C(t,Z_j(x)))\),
\(\mathcal R_+=\mathcal R(G_+,C(t,Z_j(x)))\), and
\(\mathcal R'=\mathcal R(G_+,C(t+\delta,Z_{j+1}(x)))\).
Let
\(\zeta=P_3\sqrt{\delta(1+U_*)}\) for a sufficiently large
fixed power \(P_3\) of \(\log m\).
The rectangular factor derivative bound and
\eqref{eq:sharp-discrete-oscillation} show that the
square-root lag changes by at most \(C\zeta\sqrt{T_j}\)
when the selector is moved to or from the representative,
with either \(G\) or \(G_+\) held fixed.
The state and time advances change it by
\(C\zeta(E_j+1)\) with \(G_+\) held fixed.
Use the strict clipping contraction between the two
representative comparisons, and square with a fixed
small slack. For constants \(q_1,q_2<1\),
\begin{equation}\label{eq:sharp-discrete-lag}
 \mathcal R_+\le q_1\mathcal R+C\zeta^2 T_j,\qquad
 \mathcal R'\le q_2\mathcal R+
                  C\zeta^2(T_j+E_j^2+1).
\end{equation}
Take \(A\) large enough that \(\zeta\) tends to zero faster
than any inverse-logarithmic power required below.
For \(S_j=\mathcal R+\zeta T_j\), the second inequality and
the discount recursion give
\[
 S_{j+1}\le q_3 S_j+C\zeta(1+E_j^2),\qquad q_3<1 .
\]
Since \(K'=K(1+o(\omega))\) at the end of the step, we obtain
\begin{equation}\label{eq:sharp-discrete-memory}
 K'S_{j+1}-KS_j\le-cKS_j+C\zeta K(1+E_j^2).
\end{equation}

\paragraph{One-step exponential moment.}
At an old point with state \(z=Z_j(x)\), put
\(\xi=Z_{j+1}(x)-z\).
By the derivative bounds and the uniform clipped
displacement cap, Taylor's formula gives
\[
 F(t,z+\xi)-F(t,z)\le
 v\cdot\xi+\tfrac12\xi^{\mathsf T}D_+\xi,\qquad
 D_+=D^2F(t,z)+e_m K\,\operatorname{diag}(b_i^{-2}).
\]
Here \(e_mP_1=o(\omega)\), by taking \(B\) large.
Indeed the Hessian changes along the segment by at most
\(P(1+u)\) times the relative width displacement in weighted
matrix norm, and \(1+u\le PK\). The added diagonal dominates
that change. The phase is convex, so \(D_+\) is positive
semidefinite. Advancing time at the new state gives
\[
 F(t+\delta,z+\xi)-F(t,z+\xi)\le
 -(1+\omega/2)\delta K ,
\]
by the time inequality and comparability of \(K\).

For any convex test \(\Phi\) and an even scalar attenuation
\(a\in[0,1]\) of a symmetric Gaussian vector \(g\),
\[
 \E\Phi(ag)=\tfrac12\E[\Phi(ag)+\Phi(-ag)]
 \le \tfrac12\E[\Phi(g)+\Phi(-g)]=\E\Phi(g).
\]
In fact the pair \(ag,-ag\) consists of opposite convex
combinations of \(g,-g\). With two independent Gaussian
contrasts and separate attenuations, apply this twice,
conditioning each time on the other summand. Thus the
exponential of the linear plus positive quadratic Taylor
bound has moment at most that of the uncut Gaussian sum.

Here is the quantitative Gaussian estimate. For a centered
Gaussian vector of covariance \(\Sigma\preceq\delta G_+\),
completion of the square gives, for \(p\) near one,
\[
 \frac1p\log \E e^{p(v\cdot g+g^{\mathsf T}D_+g/2)}
 =
 -\frac1{2p}\log\det(I-pA)
 +\frac p2 v^{\mathsf T}\Sigma^{1/2}
                  (I-pA)^{-1}\Sigma^{1/2}v,\qquad
 A=\Sigma^{1/2}D_+\Sigma^{1/2}.
\]
The formula is valid also for singular \(\Sigma\) by using
a standard Gaussian and the symmetric square root.
The scaled Hessian and metric bounds give
\(\|A\|\le \delta P(1+2U_*)=o(\omega)\).
Expanding the determinant and bounding the inverse show
that for \(p=1+o(\omega)\) this expression is at most
\[
 (1+o(\omega))\frac\delta2
      \tr[G_+(D_++vv^{\mathsf T})]
 \le (1+o(\omega))\delta K(1+\mathcal R_+/2)
                  +o(\omega\delta K).
\]
The last inequality uses \(B_0\preceq B=KC^{-1}\) and the
weighted bound \(e_mP_1=o(\omega)\).

Now use the exponent \(\Xi_j=F(t,Z_j)+\kappa\delta KS_j\).
Apply Hölder with \(p=1+\tau\) and its conjugate \(p'=p/(p-1)\)
to the Taylor moment and the positive memory term in
\eqref{eq:sharp-discrete-memory}. Choose
\(\zeta/\tau=o(1)\) and \(\tau=o(\omega)\).
Since \(E_j\) has a uniform conditional Gaussian tail,
\(\log \E[e^{\theta(1+E_j^2)}\mid\mathcal F_j]\le C\theta\)
for \(\theta\ge0\) sufficiently small, uniformly in the
old history \(\mathcal F_j\).
Take \(\theta=Cp'\kappa\delta\zeta K\), which is small.
The conjugate Hölder factor then costs
\(O(\kappa\delta\zeta K)=o(\omega\delta K)\) in the
logarithm. The first bound in \eqref{eq:sharp-discrete-lag}
gives \(\mathcal R_+\le C S_j\).
Choose fixed \(\kappa\) large enough that the negative
memory drift pays the lag part of the moment. The remaining
bound is at most
\[
 -(1+\omega/2)\delta K+(1+o(\omega))\delta K
                               +o(\omega\delta K)\le0.
\]
Consequently \(\E[e^{\Xi_{j+1}(x)}\mid\mathcal F_j]\le
e^{\Xi_j(x)}\) at every active point. The exponent is
frozen in full on every stopped cell.

Conditional integration over space gives a supermartingale
of total modular at least one, since the area is one and
all terms in \(\Xi_j\) are nonnegative. At zero the selector
and the metric are both the identity, and there is no
memory. The initial modular is therefore \(e^{F(0,0)}=1+o(1)\).
It follows by Markov's inequality on its nonnegative excess
that the terminal modular is \(1+o_{\mathbb P}(1)\).
On the same event, each high-stop cell carries exponent
at least \(U_*/3-O(u_0)\) throughout. Their union has area
at most \(e^{-cU_*}\). Every stored state has \(H\le2U_*\)
and lies inside the moving strip, so the horizontal level
estimate and the vertical strip bound from
Section~\ref{sec:sharp-phase} apply to the entire retained
regular field.

\subsection{Routing the rare stopped regions}
\label{sec:sc-backup}

Let \(U_*=m/(\log m)^A\), and condition on a main exploration
history for which the high-state stopped parents cover a set
\(\mathcal H\) of area \(p\le e^{-cU_*}\).
The case \(p=0\) requires no additional exploration.
The unrevealed sources in the stopped parents are independent
white noise bridges by predictable orthogonal revelation.
The old metrics, in physical balanced coordinates, satisfy
\[
 G_{Q,xx}\le (1+U_*)\operatorname{polylog}(m),\qquad
 G_{Q,yy}\le\operatorname{polylog}(m).
\]
All estimates below are conditional on that history.
We build a current for the additional contrasts and leave only
terminal bridges for the final mesh correction.

\paragraph{Putting the roots in one chart.}
Put \(a_b=1+\log(1+\log(m/U_*))\) and set
\(A_b=\diag(a_b^{-1/4},a_b^{1/4})\) and \(G_b=A_bA_b^{\mathsf T}\).
On the stopped parents, successively clip the old metric towards
\(G_b\) with relative eigenvalue ratio \(K_0\).
Concretely, clip the eigenvalues of \(G^{-1/2}G_bG^{-1/2}\)
to \([K_0^{-1},K_0]\) and conjugate back by \(G^{1/2}\).
The clipped metric is determinant one, lies between \(K_0^{-1}G\)
and \(K_0G\), and satisfies \(G_+\preceq K_0^{-1}G+G_b\),
as follows from the same scalar inequalities on the eigenvalues.
The logarithm of the larger eigenvalue of
\(G^{-1/2}G_bG^{-1/2}\) decreases by \(\log K_0\), down to zero.
Consequently every nonterminal descendant reaches \(G_b\) in
\(O(\log m)\) splits. At each step use the ordinary subdivision
and ray current of Lemmas~\ref{lem:local-geometry}
and~\ref{lem:coarse-envelope}, with the common minimum leg \(s\).
At the geometric cutoff retain the cell for the terminal correction.
Use even local clipping at \(\sqrt m/P\), where \(P\) is a
sufficiently large fixed power of \(\log m\).
The regular current has chart supremum \(C\sqrt m/P\) per step.
Iteration of the matrix inequality gives along each branch
\[
 \sum_j \sqrt{G_{j,ii}}\le C\sqrt{G_{0,ii}}
       +C(\log m)\sqrt{G_{b,ii}}.
\]
After division by \(\sqrt m\), the regular steering current
therefore has vertical supremum \(o(1)\) and horizontal
supremum at most
\((1+U_*)^{1/2}\operatorname{polylog}(m)/P\) on \(\mathcal H\).
For a field of supremum \(M_0\) on a set of area \(p\), the
quadratic Luxemburg norm is at most
\(M_0/\sqrt{\log(1+1/p)}\). This follows by substituting that
scale in its modular. The horizontal steering norm is \(o(1)\).
For the removed pieces use Lemma~\ref{lem:sc-sparse}
with the current divided by \(\sqrt m\), threshold \(\sqrt m/P\)
and box enlargement \(P^2\log m\). The removed currents
have both normalized gauges \(o(1)\) with probability \(1-o(1)\).
These corrections use the revealed contrasts. Conditional
on the steering exploration, the remaining bridges are fresh.
In \(A_b^{-1}\Omega\), the nonterminal roots have metric \(I\),
and we start their regular backup currents at zero.

\paragraph{A phase for the remaining bridges.}
Fix \(0<\alpha<1/10\), to be chosen with a clipping exponent
\(\alpha/2<\eta<1/10\). Define
\begin{align*}
 r(u)&=\frac{\sqrt u}{(1+u)^{1/2+\alpha}}
 +\frac{1}{1+\log(m/U_*)}
   \frac{(u/U_*)^\alpha}{1+(u/U_*)^\alpha}(1+u/m)^{-\alpha}\\
 &\qquad+\frac{(u/m)^\alpha}{(1+(u/m)^\alpha)^2},
 \qquad I(u)=\int_0^u r(v)\,\frac{dv}{v}.
\end{align*}
The first and third summands have bounded integrals in \(du/u\).
For the middle summand integrate separately over
\((0,U_*)\), \([U_*,m]\), and \((m,\infty)\), using the
power decay on the two tails and the prefactor on the middle
interval. Thus \(I(\infty)\le C_\alpha\). The elasticity
\(\beta=d\log r/d\log u\) lies in \([-\alpha,1/2]\)
and has uniformly bounded derivative in \(\log u\).
Each summand has these bounds, and the derivative of their
weighted mean adds the variance of the individual elasticities.
For any fixed \(C_1\), the third summand yields
\begin{equation}\label{eq:sc-backup-floor}
 r(u)\gtrsim_{C_1}m^{-\alpha},\qquad
 \frac{1+u}{mr(u)}\lesssim_{C_1}1,\qquad 1/2\le u\le C_1m .
\end{equation}

Use \(Y(y)=\sqrt{1+y^2}-1\), \(c(t)=c_0+\Lambda t\),
where \(c_0\) is a sufficiently large fixed constant.
On the strip \(Y(y)<c(t)+I(\infty)\) define \(u=H(t,x,y)\)
by the graph
\[
 Y(y)=c(t)+I(u)-r(u)(e^{x^2/u}-1).
\]
Set \(H=0\) at \(x=0\), \(Y(y)\le c(t)\).
Write \(z=x^2/u\), \(J=e^z\), \(Z=1+z\),
\(Q=r[1-\beta(J-1)+zJ]\), \(D_0=Q/(rJ)\).
The graph derivative in \(u\) is \(Q/u\), and
\(Q\asymp rJZ>0\), using \(0\le J-1\le zJ\) and
\(\beta\le1/2\). Its endpoint limits give the asserted
unique level. On these levels \(rJ\le C\) for \(\Lambda t\le1\)
and \(Q\lesssim Z\). For \(u\ge1\) the first summand gives
\(z\le C_\alpha\log(2+u)\), hence \(Z\lesssim u\).
For \(1/2\le u\le1\) it bounds \(r\) below and \(z\) above.
Take a smooth nondecreasing cutoff \(\chi\) that vanishes
below \(1/2\) and has derivative one for \(u\ge1\),
and set \(H_0=1+\chi(H)\).

For \(u\ge1\) put
\begin{gather*}
 b_x^2=uZ/Q,\qquad b_y^2=Q/Z,\qquad K=1+u/Q,\\
 D=C_D\diag(1/Z,uZ/Q^2),\qquad v=\nabla H_0,\\
 C=\sqrt{\det(D+vv^{\mathsf T})}(D+vv^{\mathsf T})^{-1}.
\end{gather*}
Extend the widths, \(K\) and the two diagonal entries of \(D\)
by smooth positive interpolation to constants below \(1/2\).
On the flat region take \(D=C_D I\). The derivative bounds
below have constant versions on the compact transition, since
\(u,z,Q^{-1}\) are bounded there. The flattened phase and
these coefficients are smooth, including across the zero-level set.

Implicit differentiation on \(u\ge1\) gives
\[
 H_x=2x/D_0,\qquad H_y=uY'/Q,\qquad H_t=-\Lambda u/Q.
\]
From \(D_0=z-\beta+(1+\beta)e^{-z}\) one obtains
\[
 |z_x|\lesssim |x|/(uZ),\qquad
 z_y=-zY'/Q,\qquad z_t=\Lambda z/Q .
\]
The first derivatives of \(D_0\) in \(z,\beta\) are bounded.
The derivatives of \(\log Q\) in \(z\) and in \(\log u\)
at fixed \(z\) are bounded by the explicit formula for \(Q\).
Differentiating the two gradient formulas, using
\(\partial_i\beta=\beta' H_i/u\) and \(|Y'|,|Y''|\le1\), gives
\[
 |H_{xx}|\lesssim Z^{-1},\qquad
 |H_{xy}|\lesssim \sqrt u/Q,\qquad
 |H_{yy}|\lesssim uZ/Q^2.
\]
For the last bound, \(|Q_y/Q|\lesssim Z/Q\) and \(Q\lesssim Z\)
handle respectively the terms with \(Q_y\) and \(Y''\).
It follows that \(D\succeq D^2 H_0\) for large \(C_D\)
and \(|b_i b_j\partial_{ij}H_0|\lesssim K\).
For \(u\ge1\), \(v^{\mathsf T}D^{-1}v\asymp u\).
The upper bound follows from the gradient formulas.
When \(z\) is bounded away from zero, \(Z H_x^2\gtrsim u\).
For small \(z\), the graph and the choice of \(c_0\) give
\(Y(y)\ge c_0-O(1)\) and \(|Y'|\gtrsim1\), so
\(Q^2H_y^2/(uZ)\gtrsim u\).
With \(b_x^2b_y^2=u\) and \(Z\lesssim u\), this gives
\(\sqrt{\det(D+vv^{\mathsf T})}\asymp K\asymp u/Q\).
In particular for \(u\ge1\)
\[
 C^{-1}=\xi\big(C_D\diag(b_i^{-2})+ww^{\mathsf T}\big),
 \qquad w=\sqrt{Q/u}\,v,\qquad \xi\asymp1,\qquad
 |b_iw_i|\lesssim\sqrt u .
\]
The cofactor formula and \(\det C=1\) give \(C_{ii}\lesssim b_i^2\),
also in the extended region.

For a determinant-one \(G\) with \(G_{ii}\le M_0 b_i^2\) write
\(\mathcal L(G)=\sqrt{\tr(GC^{-1})}\). We record the derivative
estimates used for the metric lag. On \(u\ge1\),
\begin{equation}\label{eq:sc-backup-lag}
\begin{split}
 \mathcal L^2&\lesssim_{M_0} 1+u,\qquad v^{\mathsf T}Gv\lesssim K\mathcal L^2,\\
 |b_j\partial_j\mathcal L|&\lesssim_{M_0}\sqrt{(1+u)/r},
 \qquad |\partial_t\mathcal L|\lesssim_{M_0}\Lambda\sqrt{1+u}/r .
\end{split}
\end{equation}
In the flattened region use a fixed positive \(r\) in these
bounds. To verify the derivatives, the preceding formulas
bound \(b_j\partial_j\) of \(\log u,z,\log Q,\log b_i,\log K\)
by \(C/\sqrt r\), and their time derivatives by \(C\Lambda/r\).
For instance \(b_x|z_x|\lesssim\sqrt{z/(QZ)}\),
and \(b_y|z_y|\lesssim z/\sqrt{QZ}\).
The Hessian and the explicit gradients then bound
\(|b_i b_j\partial_j w_i|\) by \(C\sqrt{u/r}\) and
\(|b_i\partial_t w_i|\) by \(C\Lambda\sqrt u/r\).
The same scaled derivative bounds as for \(\log Q\)
hold for \(\xi\), by
\(\xi^{-2}=C_D^2/u+C_D\sum_i(b_i w_i)^2/u\).
Finally \(\mathcal L\) is the Euclidean norm of the vector
with entries \(\sqrt{\xi C_D G_{ii}}/b_i\) and the
two entries of \(\sqrt\xi G^{1/2}w\).
This proves the derivative bounds without division by a small
component. The first two inequalities in
\eqref{eq:sc-backup-lag} follow from the selector expression,
using \(D+vv^{\mathsf T}\succeq vv^{\mathsf T}\) for the second.

\paragraph{Discrete evolution and modular.}
For the backup in its chart take
\(\delta=\varepsilon^2/m\) and time \(t=j\delta\), with
\(\varepsilon>0\) a sufficiently small fixed constant.
There are at most \(C_*\!m\) stages by the geometric shrink,
so fix \(T_*=C_*m\delta\). Update each active parent metric
by relative clipping towards \(C\) evaluated at its
representative. Take the ordinary split with minimum leg \(s\)
and the ray inverse of the fresh contrast.
Multiply its chart current by \(\sqrt\delta\), applying even
local clipping at \(E_0=m^{1/2-\eta}\) before the phase update.
Write \(E\) for the resulting local envelope without
\(\sqrt\delta\). The clipping gives \(E\le CE_0\) and
\(\mathbb P(E>t\mid\text{old data})\le Ce^{-ct^2}\).
The spatial state \(U\) is the sum of the retained increments,
with the chart maps applied to the increments. Stop and
freeze the entire integrand below when a cell reaches the
geometric cutoff. We also freeze globally at the first stage
whose modular exceeds a fixed multiple of \(p\).

Here are the geometric and analytic details of the induction
on active cells. Maintain the cap \(G_{ii}\le M_0 b_i^2\)
at each point of a parent, for a sufficiently large fixed \(M_0\),
and keep the widths between successive states within an
arbitrarily small fixed relative factor.
For \(x,x'\) in the same old cell, the ancestor H\"older
bounds give
\[
 |U_i(x)-U_i(x')|\le C\sqrt\delta\, b_i(t,U(x))\sqrt{\mathcal T(x)}.
 \tag{*}
\]
Here \(\mathcal T\) is a discounted sum of squared past envelopes,
with update \(\mathcal T'=\rho\mathcal T+E^2\) for some
\(\rho<1\). To see the estimate, compare the two points in the
chart of each ancestor. The fixed chart-distortion bound per
step and the sufficiently small child shrink make their
relative distance decay geometrically with the number of
intervening levels. The widths can grow by at most the
chosen small fixed factor per level. The resulting
geometrically weighted sum of envelopes is bounded by
\(C\sqrt{\mathcal T}\) by Cauchy--Schwarz, with \(\rho\)
strictly larger than the squared geometric ratio.
The retained envelopes give \(\mathcal T\lesssim E_0^2\).
The first root has \(U=0\), \(G=I\), and \(\mathcal T=0\).

Before a modular stop, each active cell has a point with
\(H\le C_1m\). Its area is at least \(cs^2\), and the
modular dominates the integral of \(e^{H_0}\), so
this follows from the modular cap \(Cp\le C\).
Use \((*)\) from this witness to the representative and every
point in the parent. The scaled derivative bounds above
give relative changes \(o(1)\) in \(1+H,b_i,K\) on these
segments. Here the spatial parameter is
\(\sqrt\delta E_0/\sqrt{r}=O(m^{\alpha/2-\eta})=o(1)\)
by \eqref{eq:sc-backup-floor}. For rigor, impose factor-two
bounds on the coefficients and \(1+H\) along the segment,
apply the derivative bounds up to their first possible
failure using the floor up to \(2C_1m+2\), then continue.
The wall margin in the \(Y\) variable is
\[
 I(\infty)-I(u)+r(u)(J-1)\gtrsim rJ\asymp b_y^2 .
\]
For the first term integrate on \([u,2u]\), where bounded
elasticity keeps \(r\) comparable. The displacement in the
vertical direction is \(O(\sqrt\delta b_y E_0)=o(b_y^2)\)
since \(b_y\gtrsim\sqrt r\). Thus the segment stays in the
strip, with a fixed positive margin used in the flattened region.
The time step costs \(O(\Lambda\delta/r)=o(1)\) in relative
changes and widens the strip. This proves the estimates at
all old active points without presupposing a pointwise
upper bound on \(H\).

Representative comparison and
\(G_+\preceq K_0^{-1}G+C\) now give a strict margin in the
diagonal cap for the new metric. The fresh displacement
is \(O(\sqrt\delta b_i E)\), so the same continuation
argument closes the diagonal cap at the new state and
bounds the relative width changes. It also bounds the
one-step overshoot of \(H\) by \(C m\).
All chart distortions and physical cell axes have polynomial
bounds. This follows from the formulas for \(b_i\),
\eqref{eq:sc-backup-floor} and the determinant-one caps,
including the low-level extension. Thus the terminal cells
have physical diameter \(h\operatorname{poly}(m)\) by the
geometric cutoff. These bounds include the \(O(\log m)\)
steering generations.

To check the lag update, first evaluate both metrics at
the representative. If \(e^d\ge1\) is the larger relative
eigenvalue, then the squared lag is \(e^d+e^{-d}\).
Clipping decreases \(d\) by \(\log K_0\), down to zero.
Consequently the lag norm contracts by \(K_0^{-1/2}\)
up to a bounded additive constant. Changes from the
representative and on the fresh spatial and time segments
have, by \eqref{eq:sc-backup-lag}, total bound
\[
 C\sqrt{\delta(1+u)/r}(\sqrt{\mathcal T}+E)
       + C\Lambda\delta\sqrt{1+u}/r
 \le C\varepsilon(\sqrt{\mathcal T}+E)+C\Lambda\varepsilon^2 .
\]
For the last inequality use \eqref{eq:sc-backup-floor} and
its flat-region version. Take
\(S=\mathcal L^2+\mathcal T\). Squaring with slack in the
clipping contraction and using the discount update gives
\(S'\le q_0 S+C(1+E^2)\) with \(q_0<1\), for small
\(\varepsilon\) with \(\Lambda\varepsilon^2\le1\).
The relative change of \(K\) then gives
\[
 K'S'-KS\le-cKS+CK(1+E^2).
 \tag{**}
\]

Use the exponent
\[
 \Xi=H_0(t,U)+C_0(T_*-t)+\kappa\delta KS .
\]
For the fresh step move first in space, then in time.
The Taylor remainder in space is at most \(C\delta K E^2\)
by the Hessian bound and comparability along the segment.
For the predictable linear tilt the uncut conditional
Gaussian variance is at most \(C\delta K\mathcal L^2\),
by \eqref{eq:sc-backup-lag}, the coarse covariance bound
and the relative bound between \(G_+\) and \(G\).
Even scalar attenuation gives the same upper bound for the
linear exponential moment by averaging the two Gaussian signs.
In time the derivative of \(H_0\) equals \(-\Lambda u/Q\)
above the transition, and is nonpositive in the transition.
The scalar clock covers the transition and flat region for
\(C_0\ge C\Lambda\). By comparability the combined time
contribution is at most \(-c\Lambda\delta K\).
Together with \((**)\), conditional Cauchy--Schwarz
bounds the log moment of \(\Xi'-\Xi\) by
\[
 -c\Lambda\delta K-c\kappa\delta KS
       +C\delta K\mathcal L^2+C_\kappa\delta K .
\]
Indeed, the envelope tail gives
\(\log\E[\exp(aE^2)\mid\text{old data}]\le Ca\) for all
sufficiently small fixed \(a\ge0\), and
\(\delta K\lesssim\varepsilon^2\) by the floor bound.
Choose \(\kappa\) to pay the term with \(\mathcal L^2\)
and then \(\Lambda,C_0\) to pay the bounded term.
Choose \(\varepsilon\) last, so the envelope moment bound
holds and \(\Lambda T_*\le1\). All constants in the
phase estimates were uniform for this time interval.
The resulting log moment is nonpositive.

Spatial integration with the prescribed freezing proves
that \(\int_{\mathcal H}e^\Xi\) is a nonnegative
supermartingale. Initially \(H_0=1\), the scalar clock is
bounded, and \(K,\mathcal L\) are bounded at \(U=0,G=I\).
The initial integral is therefore at most \(Cp\), with
an absolute bound after shrinking \(\varepsilon\).
Choose the modular stop at a sufficiently large fixed
multiple of \(p\). The maximal inequality gives conditional
probability bounded below of completing without this stop.
On this event \(H\le Cm\) by the area witness and overshoot
estimate, \(\int_{\mathcal H}e^H\le Cp\), and the normalized
backup vertical component is bounded by the strip width \(C\).

\paragraph{Cost and exact divergence.}
On \(1/2\le u\le Cm\) the middle profile term yields
\[
 \log(C/r(u))\le C'+\log(1+\log(m/U_*))
       +\alpha\log_+(U_*/u).
\]
Indeed \((1+u/m)^{-\alpha}\) is bounded below here and
\(v/(1+v)\ge\tfrac12\min(v,1)\) for \(v=(u/U_*)^\alpha\).
Using \(rJ\le C\) and \(u\log_+(U_*/u)\le U_*/e\), we get
\[
 x^2\lesssim a_b(1+u)+U_* .
\]
This bound holds also for \(u<1/2\), because
\(r(u)\gtrsim\sqrt u\) there bounds \(u\log(C/r(u))\).
At scale \(C_2\sqrt{a_b}\), with \(C_2\) sufficiently large,
the horizontal exponential integral on \(\mathcal H\) is at
most \(C p\exp(CU_*/(C_2^2a_b))=o(1)\), by
\(\int_{\mathcal H}e^H\le Cp\) and \(p\le e^{-cU_*}\).
Outside \(\mathcal H\) the backup current is zero.
Thus its normalized horizontal \(\psi_2\) cost is
\(O(\sqrt{a_b})\).

Lemma~\ref{lem:sc-sparse} applied in the backup chart
to the removed pieces, with \(E_0=m^{1/2-\eta}\) and
box enlargement \(C m^{2\eta}\), gives normalized hard
vertical cost \(O(1)\) and horizontal \(\psi_2\) cost
\(O(m^{5\eta-1/2})=o(1)\). These estimates and the steering
repairs hold on events of conditional probability \(1-o(1)\),
uniformly before the modular stop, so intersecting with them
still leaves a fixed positive conditional success probability.
The sum of the retained and repaired currents has divergence
equal to the sum of all revealed contrasts, because each
repair replaces exactly the divergence removed by clipping.
Undo the factor \(\sqrt\delta\) and push the backup forward
by \(A_b\). Its physical costs in both directions are
\(O(\sqrt m\,a_b^{1/4})\), while the steering costs are
\(o(\sqrt m)\). Since \(a_b=O(1+\log\log\log m)
=o(\log\log m)\), these are lower order at the sharp scale.

\subsection{Assembly and exact repair}
\label{sec:sharp-gauss-upper}

Write $m=\log(1/h)$, $a=\log\log m$, and $\lambda=m/(4\pi)$.
The target is
\[
 \mathcal F(P_hW^0)\le (1+o(1))\sqrt{m/(2\pi)}\,a^{1/4}
\]
with probability bounded away from zero. Here $\mathcal F$ is the
continuum gauge with zero normal flux. Uniform Lipschitz concentration
then gives the same asymptotic upper bound in expectation and in
probability. Indeed affine interpolation of a discrete cut-capacity
solution proves $\mathcal F(h^{-1}g)\le C\|g\|_2$ for a centered
fine-grid cell vector $g$, so the Lipschitz constant with respect to
the standard Gaussian cell coefficients is uniform.

Set $L\sim(\log m)^2$, $R=e^L$ an integer, and $s=hm^{C_s}$.
Choose a determinant-one balancing map
$A_0=\diag(b,b^{-1})$ with $b^2$ an integer and
$b=(1+o(1))\sqrt2\,a^{1/4}$. In $A_0^{-1}\Omega$ tile the area-one
rectangle by $b^2$ squares of side $1/b$. Reveal their coarse
means. In physical coordinates their density is constant on
horizontal strips, and its vertical primitive has supremum
\(O_{\mathbb P}(1)\). In fact the successive strip masses before
centering are independent Gaussians of variance \(1/b^2\).
If \(S_j\) is their partial sum, the primitive at strip boundaries
is \(S_j-(j/b^2)S_{b^2}\). The \(L^2\) martingale maximal
inequality bounds its supremum in \(L^2\).
The regular evolution starts at zero, independently of these
coarse means.

At a stage, before revealing new contrasts, stop a parent on
the high-state criterion $H(t,U(x_Q))\ge U_*/2$ with
$U_*=m/(\log m)^A$. The within-parent phase comparison of the
discrete stability estimates justifies this representative
test. On each other parent clip its old metric towards the
sharp selector with the fixed relative ratio $K_0$. In the new
chart take an ordinary fixed-cap Whitney split with minimum
side $Rs$. On every nonterminal base square $S$ refine further
by the uniform $R\times R$ grid. Both partitions are chosen
before new reveals. The coarse child-average contrast and
the centered fine contrasts inside each base square are
orthogonal Gaussian projections, with independent fine
contrasts in different base squares. Use the ray inverse of Section~\ref{sec:sharp-local-inverse} for
the coarse contrast and the inverse with radial principal part for
each fine contrast. All fine children are squares
in the current metric chart. Every child of a coarse split receives the coarse update, including
the terminal base cells. These terminal cells receive no fine contrast
and freeze after the coarse update. A parent too small for a thick
split freezes directly and its unrevealed bridge goes to the terminal
repair.

Clip the coarse increment by the even envelope rule in
Lemma~\ref{lem:sc-sparse} at $\sqrt m/P$, with $P$ a sufficiently large power of
$\log m$. For the fine contrast clip a whole base square if any
fine-child envelope, with amplitude and oscillation divided by
$\sqrt L$, exceeds $\sqrt m/P$. On an accepted base square retain
the entire fine current. It has zero normal flux on that square,
so rejecting it removes exactly its fine source.
Each fine rejection has conditional probability at most
$\exp(-c m/P^2)$ by the Gaussian envelope bounds and the union
over $R^2$ children. These events are conditionally independent
over base squares, including after the coarse reveal. The
coarse and fine attenuation rules depend separately on their
own innovations and are even in those innovations.

The chart covariance of the uncut sum is bounded above by
$(L/(4\pi)+O(\sqrt L))I$ by the isotropic-principal lemma and
the $O(1)$ coarse covariance. Divide the physical chart increment
by $\sqrt\lambda$ in the phase evolution, so the step clock is
$\delta=(L+O(\sqrt L))/m$. Every nonterminal branch shrinks by
a factor $O(R^{-1})$ in intrinsic leg per step, since the relative
metric distortion is fixed. The depth is at most
$(m+O(\log m))/(L-O(1))$. The sharp phase can therefore run over
a deterministic horizon $T=1+o(1)$. Its discrete stability gives
terminal horizontal $\psi_2$ norm at most $1+o(1)$ and vertical
supremum at most $(2+o(1))\sqrt a$ in normalized unbalanced
coordinates. The former follows from $x^2\le(1+o(1))H+o(1)$
on the active levels and their one-step overshoots, and the
terminal modular $1+o_{\mathbb P}(1)$. The latter follows from
the moving strip, whose final radius is
$L_0+2(1+o(1))\sqrt a\,T+o(1)$. These bounds include the stored
regular values on the high-stop cells. Their union has area at
most $\exp(-cU_*)$. The balancing map and the factor
$\sqrt\lambda$ give exactly
$(1+o(1))\sqrt{m/(2\pi)}a^{1/4}$ in both physical gauges.

\paragraph{The removed fine sources.}
Use physical determinant-one repair maps
$D_k=\diag(k,k^{-1})$, with a sufficiently large polylogarithmic
$k$ to be fixed below. For each candidate base square, in
$D_k^{-1}\Omega$ choose a predictably enlarged axis square
containing its transformed support. Its side is at most a
polynomial in $m$ times the intrinsic base leg. If it cannot
be fit in the domain, the base leg is inverse-polynomial and
the conditional rejection union bound rules out rejection
of any such square with probability $1-o(1)$.
For a fitting candidate subdivide the repair square dyadically down
to width $g$, comparable to its fine leg times a fixed inverse
polynomial in $m$. The logarithm of the number of mesh cells per side is
$L+O(\log m)=o(\sqrt m)$.

For every candidate fine contrast $f$, even before testing
its rejection, the Gaussian vector of its projected repair-cell
masses divided by $g$ has covariance at most the identity.
The coordinate tests are disjoint cell indicators divided by
$g$, composed with the determinant-one map and projected onto
the fine contrast space in its original base square. Orthogonal
projection is a contraction. The sum of the masses is zero
in the repair square. The capacity estimate of Lemma~\ref{lem:sc-box-gauss} and Gaussian
concentration give common repair load $O(\sqrt m)$ with
conditional failure probability $e^{-C m}$ for arbitrarily
large fixed $C$. Union over all attempted candidates.
The residual $(I-P_g)f$ is mean zero on each repair cell and
costs at most $Cg\|f\|_\infty$ by coordinate primitives.
The Gaussian union bound gives $\|f\|_\infty\le C\sqrt m$
divided by the fine leg, so this cost is $O(\sqrt m)$.
Thus on the joint good event, the exact repair per rejected
square has physical horizontal supremum $O(k\sqrt m)$ and
vertical supremum $O(\sqrt m/k)$ by pushing forward $D_k$.

At any point and stage, polynomially many candidate repair
boxes cover it, by disjoint area packing in dyadic intrinsic
size classes and polynomial distortion. Conditional independence
of fine rejections gives the exponential overlap bound by
expanding the one-stage counting moment at tilt $c m/P^2$.
Iterating over the stages and union over a fine net after
padding boxes gives a uniform overlap $O(P^2)$. The expected
sum of rejected box areas is polynomial times
$e^{-c m/P^2}$, so their union has area at most
$e^{-c' m/P^2}$ with high probability. Choose $k/P^2\to\infty$
polylogarithmically. The vertical repair is $o(\sqrt m)$.
The horizontal supremum is polylogarithmic times $\sqrt m$
on this sparse support, hence its global $\psi_2$ norm is
polylogarithmic and is also $o(\sqrt m)$.
Lemma~\ref{lem:sc-sparse}, applied to the unscaled coarse
current with threshold \(\sqrt m/P\) and repair stretch
\(k\gg P^2\) polynomial in \(\log m\), gives
vertical cost \(o(\sqrt m)\) and horizontal Luxemburg cost
\(o(\sqrt m)\) for the removed coarse sources. Treat the two types
separately, so no joint coarse-fine rejection estimate is needed.

\paragraph{Completion.}
Use the shifted-tail backup on the stopped union, conditional
on the full main history. It succeeds at fixed positive
conditional probability, with both physical budgets
$o(\sqrt m a^{1/4})$. Its initial steering increments and their
sparse repair have smaller order as well. Steering and backup
use the same minimum side $s$ as the main exploration. The
terminal leaves of the combined exploration have physical
diameter at most \(hR\operatorname{poly}(m)\).
Before the last residual projection, the revealed density satisfies
\(\|f_{\rm leaf}\|_\infty\le C m^{3/2}/s\) with high probability.
Each contrast has density standard deviation at most \(C/s\)
on its new cells, all of which have area \(\gtrsim s^2\).
A Gaussian union bound over \(O(m)\) generations with at most
\(Cs^{-2}\) cells per generation gives this estimate by summing
the contrasts along the branches. The initial root-square means
satisfy the same bound. Thus the deterministic correction
\((P_h-I)f_{\rm leaf}\) costs \(O(\sqrt m)\) by coordinate
primitives on the final grid. For the remaining bridges choose
an $h$-aligned dyadic grouping grid of side $D$ comparable to
the terminal diameter bound and larger than $h$ plus that bound.
Assign each leaf to the grouping square containing a representative
of the leaf. Enclose the square and its immediate neighbors in
a square of side \(4D\) with sides on the \(h\)-grid. At the boundary
slide that square inside \(\Omega\). Each enlarged square is within
\(O(D)\) of the original one, so their overlap is bounded. They
contain the final mesh projections of every assigned bridge. Each projected bridge has zero mass.
Conditional on the exploration, the normalized mesh masses of the
sum of the assigned bridges have covariance at most the identity,
by orthogonal projection onto their unrevealed Gaussian subspaces.
Lemma~\ref{lem:sc-box-gauss} thus gives mean capacity
$O(\log(D/h))=O(L+\log m)=o(\sqrt m)$ in each box.
Uniform Gaussian concentration and a union bound give total
repair cost $O(\sqrt m)$. All conditional bridges
are valid since the geometry, stopping, and discarded increments
have used only already revealed masses. Sum the coarse means,
main regular current, exact removed-source repairs, backup and
terminal currents. Their divergence is exactly $P_hW^0$.



The total current has both gauges at most
\((1+o(1))\sqrt{m/(2\pi)}\,a^{1/4}\) on an event of probability
bounded below. By the uniform Gaussian concentration of the gauge,
its mean is at most this bound plus \(O(1)\), and the same bound
with a relative \(o(1)\) holds in probability. Together with the
lower bound of Section~\ref{sec:sharp-constant-lower}, this proves
Theorem~\ref{thm:gaussian}.

\section{An asymptotically exact Gaussian reduction}
\label{sec:exact-transfer}

We first compare the point cost with the flux gauge of an empirical
histogram. The upper inequality needs only local count estimates.
Once we couple that histogram to white noise, the Gaussian upper
bound gives the order estimate required for the reverse inequality.

\begin{theorem}\label{thm:exact-transfer}
Put \(m=\log N\), \(s_N=\sqrt{m/N}\) and
\(\rho_N=(\log\log m)^{1/4}\).
For \(\ell\) dyadic with
\(\ell\asymp N^{-1/2}m^{10}\), the endpoint cost and white
noise can be coupled so that
\begin{equation}\label{eq:exact-coupling}
 M_N=\sqrt{2/N}\,\mathcal F(P_\ell W^0)+O_{\mathbb P}(s_N).
\end{equation}
In particular, if \(\ell=2^{-j}\), then
\begin{equation}\label{eq:exact-mean-reduction}
 \frac{M_N}{s_N\rho_N}-
 \frac{\mathbb E\mathcal F(P_{2^{-j}}W^0)}
 {\sqrt{j\log 2}\,(\log\log j)^{1/4}}
 \ \xrightarrow{\mathbb P}\ 0 .
\end{equation}
The same statement holds after exchanging the two directions.
\end{theorem}

The proof uses the conditional count estimate of
Lemma~\ref{lem:sc-box-emp} and the Gaussian upper bound proved above.
All costs use the convention in Theorem~\ref{thm:main}.

\subsection{From point measures to a smooth flux gauge}

Let \(\mu_N^+,\mu_N^-\) be the two empirical probability
measures and choose a dyadic \(h\asymp C_h s_N\) with a
sufficiently large fixed \(C_h\). Let \(S_\ell\) smooth by
reflection on the square. Explicitly, extend the source to the
period-two torus by even reflection in each coordinate, then
convolve with a nonnegative, reflection-symmetric kernel
\(K_\ell\) of integral one, supported in a ball of radius
\(C\ell\). Take \(\|K_\ell\|_\infty\le C\ell^{-2}\) and
\(\|\nabla K_\ell\|_\infty\le C\ell^{-3}\). For vector currents
reflect the normal component oddly and the tangential component
evenly. This extension commutes with divergence, and the
smoothed current has zero normal trace on the square.

\paragraph{Uniformization.}
Each sample couples to its \(P_\ell\) histogram with
coordinatewise maximum displacement \(O(s_N)\) with probability
tending to one. First spread each point uniformly in its
\(h\)-cell. The \(P_h\) histogram is bounded below by a positive
absolute constant on the uniform occupancy event. Within each
\(\ell\)-square, compare the fine counts to their mean over the
square. Conditional on its total, these are multinomial counts
with uniform cell labels. Lemma~\ref{lem:sc-box-emp} gives
a signed flow with both face loads \(O(Nh^2)\), uniformly over
the \(\ell\)-squares with probability tending to one.
The estimate applies since the boxes have polynomially many
\(h\)-cells per side, their conditional cell means are
\(O(Nh^2)\), and the conditional failure bound has arbitrarily
large exponential rate in \(m\) after increasing the load bound.
The lower bound on \(P_h\mu_N^\pm\) also bounds
\(P_\ell\mu_N^\pm\) below. The face interpolation gives a
zero-normal flux of supremum \(O(h)\) for the difference of the
two histograms. Linear interpolation of the densities and
velocity equal to flux divided by density then give the desired
coupling. For a measurable flux this velocity argument follows
by first smoothing with the same reflection convention at a
scale tending to zero and taking a weak limit of couplings.

Chernoff's inequality on the coarse grid also gives
\[
 \|P_\ell\mu_N^\pm-1\|_\infty=O(s_N/\ell)
\]
with probability tending to one. The histogram \(p\) couples
to \(S_\ell p\) at maximum displacement \(O(s_N)\). The
evenly extended density has translation flux
\[
 J(z)=\int K_\ell(v)\,v
                 \int_0^1(p(z-tv)-1)\,dt\,dv
\]
with divergence \(p-S_\ell p\) and supremum at most
\(C\ell\|p-1\|_\infty\). The parity of \(J\) gives the
boundary condition, and the linear-density velocity again
gives the coupling. Finally reflection smoothing is
contractive for the coordinatewise maximum transport distance.
Add the same random displacement and fold both positions
back into the square with the coordinatewise period-two
reflection map, which is \(1\)-Lipschitz in each coordinate.
By gluing, each sample couples to its \(S_\ell\) density at
maximum displacement \(O(s_N)\).

For use in the other direction, take a coupling whose starting
marginal is one of the samples and whose maximum displacement is
\(D=o(\ell)\). Spread its straight paths with the reflected
kernel. The time-\(t\) occupation density \(r_t\) then differs
uniformly from the smoothed starting density by \(O(D/\ell)\).
Indeed, only starting atoms within distance \(C\ell\) of the
evaluation point or its reflections can contribute. These
have total mass \(O(\ell^2)\) by the coarse occupancy cap.
The derivative bound for the kernel now gives the estimate.
Apply this also to the point-to-\(P_\ell\) coupling with
\(D=O(s_N)\). It follows that
\begin{equation}\label{eq:smooth-densities}
 \|S_\ell\mu_N^\pm-1\|_\infty=O(s_N/\ell).
\end{equation}

\paragraph{Sharp upper transfer.}
For densities \(p,q\) each bounded below by \(1-\epsilon\),
a zero-normal flux of divergence \(p-q\) and gauge \(r\)
gives a coupling of gauge at most \(r/(1-\epsilon)\).
For the velocity \(J/\rho_t\) along
\(\rho_t=(1-t)p+tq\), convexity gives
\[
 \int \!\left(e^{((1-\epsilon)\Delta_1/r)^2}-1\right)d\pi
 \le \int_0^1\!\int \rho_t
       \left(e^{((1-\epsilon)J_1/(r\rho_t))^2}-1\right) dx\,dt
 \le (1-\epsilon)\int\!\left(e^{(J_1/r)^2}-1\right) dx .
\]
The hard bound follows by integrating the absolute vertical
velocity on each path. Smooth approximation as above proves
this for general feasible fluxes. Use \eqref{eq:smooth-densities}
and glue to the point measures with additive cost \(O(s_N)\).
The resulting fractional point coupling has the hard
displacement cap as a support constraint. Multiply its matrix by
\(N\), giving row and column sums one. Any row subset has at least
as many positive-support neighbors as rows by the two marginal
conditions, so Hall's theorem gives a supported permutation.
Subtract the smallest entry on that permutation, rescale the
remaining common marginal, and repeat. Each subtraction either
finishes the decomposition or removes a positive entry. The
horizontal modular is linear in the coupling, so one of the
supported permutations has modular at most that of the coupling.
It follows that
\[
 M_N\le (1+O_{\mathbb P}(s_N/\ell))
   \mathcal F(S_\ell(\mu_N^+-\mu_N^-))+O_{\mathbb P}(s_N).
\]
The smoothed flux gauge may be replaced by the histogram one
with additive error \(O_{\mathbb P}(s_N)\). For each color
spread the paths of its point-to-\(P_\ell\) coupling.
Its smoothed occupation density is uniformly bounded, and
the resulting flux has both suprema \(O(s_N)\). It connects
\(S_\ell\mu_N^\pm\) to \(S_\ell P_\ell\mu_N^\pm\).
The translation flux then connects the latter to
\(P_\ell\mu_N^\pm\) at the same cost.
By subadditivity and symmetry of \(\mathcal F\), the two
empirical flux gauges differ by \(O_{\mathbb P}(s_N)\).
\subsection{Sharp coupling of the histogram counts}

Poissonize the coarse cell counts of each color, so they
are independent Poisson variables of mean \(\lambda=N\ell^2\).
Each Poisson count can be coupled by its quantile to a
standard normal \(Z\) so that
\begin{equation}\label{eq:pois-quantile}
 |P-\lambda-\sqrt\lambda Z|\le C m^{3/2}
 \qquad\text{on } |Z|\le 4\sqrt m .
\end{equation}
Here only \(\lambda\asymp m^{20}\) is needed.
For a proof put \(R=4\sqrt m\). Stirling's formula, expanding
the Poisson rate function at \(\lambda\), compares each mass
at \(|k-\lambda|\le8R\sqrt\lambda\) to its normal cell mass
with relative error \(O((1+R^3)/\sqrt\lambda)\).
The normal cell uses endpoints
\((k-\lambda\pm 1/2)/\sqrt\lambda\).
The Poisson and normal tails beyond this range are negligible
relative to the relative-error tolerance times either retained tail
at a threshold in \([-2R,2R]\), by
Chernoff's inequality. Summing the mass estimate gives the
same relative error for both distribution tails in that
threshold range. The lower normal tail has logarithmic
derivative bounded away from zero on \((-\infty,1]\), and
the upper tail has the analogous bound on \([-1,\infty)\).
Hence thresholds \(Z\pm C(1+R^3)/\sqrt\lambda\), with the
rounding error included, bracket the Poisson quantile.
This proves \eqref{eq:pois-quantile}.

Couple the cells independently and truncate each coupling
error to zero on \(|Z|>4\sqrt m\). With probability tending
to one this changes none of the cells. Center the errors
by removing the spatial mean of the histogram.
For any dyadic box of side \(d\ge\ell\), the sum of errors
centered by their expectations is \(O(m^2 d/\ell)\) uniformly
over the boxes, with probability tending to one, by
Hoeffding and a union bound. Spatial centering costs only
a constant factor in this estimate.
The histogram residual divided by \(N\) therefore has a
flux with both suprema \(O(m^3/(N\ell))\).
In detail, telescope the averages over the dyadic levels.
For a mean-zero child-average difference in a box of side
\(d\), row primitives after subtracting row means, followed
by a vertical primitive of those row means, have suprema
at most \(Cd\) times the maximum child density.
Each level costs \(O(m^2/(N\ell))\) and there are \(O(m)\)
levels. This residual is \(o(s_N)\).

To restore fixed sample size, condition on the Poisson
total \(K\) and put the Poissonized points in independent
uniform order, then append independent uniform points as
needed. Conditional on \(K\) this sequence is iid uniform,
so taking its first \(N\) points gives the required law.
The difference has \(|K-N|\le\sqrt N m\) extra or removed
uniform labels with probability tending to one.
If there are \(k\) such labels, the centered count in each
dyadic box of side \(d\) is
\(O(\sqrt{k d^2m}+m)\) uniformly by Bernstein.
The same dyadic primitives for this centered correction cost
at most
\[
 C m\sqrt{km}/N+C m/(N\ell)=o(s_N).
\]
All comparisons remove the spatial mean, which cancels in
the difference of the two fixed-size samples.

In continuum units the Gaussian count term for one color
is \(N^{-1/2}P_\ell W^0\), since the indicator of a cell has
\(L^2\) norm \(\ell\) and the raw count standard deviation is
\(\sqrt N\ell\). Independence of the two colors gives
\(\sqrt{2/N}P_\ell W^0\) for the difference. Thus the empirical
and Gaussian histogram gauges differ by \(o_{\mathbb P}(s_N)\)
after this factor. In particular,
\begin{equation}\label{eq:hist-gauss}
 \mathcal F(P_\ell(\mu_N^+-\mu_N^-))
 =\sqrt{2/N}\,\mathcal F(P_\ell W^0)+o_{\mathbb P}(s_N).
\end{equation}
The Gaussian upper bound and the upper transfer now yield
\(M_N=O_{\mathbb P}(s_N\rho_N)\).


\subsection{Reverse inequality and the normalization}
Spread the paths of an endpoint-optimal permutation with the
same kernel, reflecting each path and its direction as above.
Its maximum horizontal displacement is at most
\(M_N\sqrt{\log(2N)}\) by its exponential modular, and its
maximum vertical displacement is at most \(M_N\).
Thus \(D=O_{\mathbb P}(m\rho_N/\sqrt N)=o_{\mathbb P}(\ell)\).
The occupation error \(O(D/\ell+s_N/\ell)\), multiplied
by \(M_N\), is \(o_{\mathbb P}(s_N)\).
By the occupation estimate,
\(\int_0^1 r_t\,dt\le1+o_{\mathbb P}(1)\) uniformly.
The smoothed path current has divergence
\(S_\ell(\mu_N^+-\mu_N^-)\). Pointwise Jensen, with zero
filler so that the weights have total mass one, bounds its
horizontal Luxemburg norm by \((1+o_{\mathbb P}(1))M_N\).
Indeed, for an occupation density bound \(1+\epsilon\) one
uses the scale \((1+\epsilon)M_N\). Integrating the Jensen
bound in space counts each path with its original weight
\(1/N\), since the reflected kernel preserves that mass on the
square. The vertical supremum has the same factor by direct
averaging. Hence
\[
 \mathcal F(S_\ell(\mu_N^+-\mu_N^-))
               \le M_N+o_{\mathbb P}(s_N) .
\]


By the histogram comparison, the last inequality also holds for
\(\mathcal F(P_\ell(\mu_N^+-\mu_N^-))\) with an \(O_{\mathbb P}(s_N)\)
error. The multiplicative error in the upper transfer is
\(O_{\mathbb P}(s_N/\ell)\) times a gauge of order
\(O_{\mathbb P}(s_N\rho_N)\), hence it is \(o_{\mathbb P}(s_N)\).
Together with \eqref{eq:hist-gauss} this proves \eqref{eq:exact-coupling}.

Write \(\ell=2^{-j}\), so
\(j\log 2=\frac12\log N-10\log\log N+O(1)\).
The uniform Gaussian Lipschitz estimate for \(\mathcal F\), obtained
by applying the grid cut bound \eqref{eq:sc-box-cut} and affine
interpolation, gives fluctuations
\(O_{\mathbb P}(1)\) around the expected flux gauge.
Since \(\log\log j\sim\log\log\log N\), this proves
\eqref{eq:exact-mean-reduction}. Theorem~\ref{thm:gaussian} now gives
\[
 M_N=\left(\frac{1}{\sqrt{2\pi}}+o_{\mathbb P}(1)\right)
         \sqrt{\frac{\log N}{N}}\,(\log\log\log N)^{1/4}.
\]
The factor \(\sqrt{2/N}\) from the difference of the samples and the
factor \(\sqrt{\frac12\log N}\) from the cutoff combine to
\(\sqrt{\log N/N}\). Reflection in the diagonal of the square preserves
the joint sample law and exchanges the coordinate costs, giving the
other endpoint.

\begin{thebibliography}{9}
\bibitem{umc:repository}
\url{https://github.com/grottofrancesco/UltimateMatchingTheorem}.
\end{thebibliography}

\end{document}
